\documentclass[11pt,a4paper]{article}
\usepackage[T1]{fontenc}
\usepackage[utf8]{inputenc}
\usepackage{lmodern}
\usepackage{microtype}
\usepackage[margin=27mm,headheight=14pt]{geometry}
\usepackage{amsmath,amssymb,amsthm,mathtools}
\usepackage{booktabs,array,longtable}
\usepackage{xcolor}
\usepackage{enumitem}
\usepackage{fancyhdr}
\usepackage{listings}
\usepackage{hyperref}
\usepackage{bookmark}
\definecolor{ink}{HTML}{16324F}
\definecolor{accent}{HTML}{19647E}
\definecolor{pale}{HTML}{EFF5F8}
\hypersetup{colorlinks=true,linkcolor=ink,citecolor=accent,urlcolor=accent,
 pdftitle={Small-prime congruences for OEIS A348410},
 pdfauthor={ChatGPT},pdfsubject={The cubic prime-power tower, an exact binary defect formula, and a sharp denominator-12 theorem},
 pdfkeywords={OEIS A348410, supercongruence, prime powers, Mobius transform, generating functions, reciprocal squares}}
\pagestyle{fancy}
\fancyhf{}
\fancyhead[L]{\small\textsc{Small-prime congruences for A348410}}
\fancyhead[R]{\small Research article}
\fancyfoot[C]{\thepage}
\renewcommand{\headrulewidth}{0.35pt}
\setlength{\parskip}{4pt}
\setlength{\parindent}{1.2em}
\setlength{\emergencystretch}{2em}
\setcounter{tocdepth}{2}
\numberwithin{equation}{section}
\newtheorem{theorem}{Theorem}[section]
\newtheorem{lemma}[theorem]{Lemma}
\newtheorem{proposition}[theorem]{Proposition}
\newtheorem{corollary}[theorem]{Corollary}
\theoremstyle{definition}
\newtheorem{definition}[theorem]{Definition}
\newtheorem{problem}[theorem]{Problem}
\newtheorem{example}[theorem]{Example}
\theoremstyle{remark}
\newtheorem{remark}[theorem]{Remark}
\newtheorem{conjecture}[theorem]{Conjecture}
\newcommand{\Z}{\mathbb Z}
\newcommand{\Q}{\mathbb Q}
\newcommand{\R}{\mathbb R}
\newcommand{\CT}{\operatorname{CT}}
\newcommand{\vp}{v_p}
\newcommand{\eps}{\varepsilon}
\newcommand{\coeff}[2]{[#1^{#2}]}
\newcommand{\A}{A_{\alpha,\beta}}
\newcommand{\B}{B_{\alpha,\beta}}
\newcommand{\dd}{\,\mathrm d}
\DeclareMathOperator{\Res}{Res}
\lstset{basicstyle=\ttfamily\small,columns=fullflexible,breaklines=true,
 frame=single,rulecolor=\color{ink!25},backgroundcolor=\color{pale},
 xleftmargin=0.5em,xrightmargin=0.5em,aboveskip=1em,belowskip=1em}

\begin{document}
\begin{titlepage}
\thispagestyle{empty}
\vspace*{11mm}
{\color{accent}\large\scshape Combinatorics and arithmetic of integer sequences\par}
\vspace{10mm}
{\color{ink}\LARGE\bfseries Small-prime congruences\par
for OEIS A348410\par}
\vspace{6mm}
{\Large The cubic prime-power tower, an exact binary\par
defect formula, and a sharp denominator-12 theorem\par}
\vspace{10mm}
{\large Research article prepared with ChatGPT\par}
{19 September 2026\par}
\vspace{9mm}
\hrule
\vspace{5mm}
\begin{abstract}
Let $a(n)=[x^n]((1-x)(1-x^2))^{-n}$, the basket sequence OEIS A348410,
and more generally
$\A(n)=[x^n](1-x)^{-\alpha n}(1+x)^{-\beta n}$ for arbitrary integer
parameters. An OEIS comment of Peter Bala conjectures the cubic congruence
$a(mp^r)\equiv a(mp^{r-1})\pmod{p^{3r}}$ for primes $p\ge5$.

We prove two theorems. First, for every odd prime $p$ and every $N$
divisible by $p$, the difference $\A(N)-\A(N/p)$ is divisible by
$p^{3v_p(N)-\eps_p}$, where $\eps_3=1$ and $\eps_p=0$ for $p\ge5$; no
coprimality restriction on the multiplier is needed. This contains Bala's
assertion. A prior public proof candidate for the same odd-prime statement
was located during the source audit and is credited below; no priority is
claimed for it. Second, at the prime left untreated in that source we prove
an explicit first binary defect congruence modulo $2^{3r-2}$, which gives a
sharp universal modulus $2^{3r-3}$, improved to $2^{3r-2}$ when $\alpha\beta$
is even.

Together these yield $12n^{-3}\sum_{d\mid n}\mu(d)a(n/d)\in\Z_{\ge0}$ for
every $n\ge1$, with $12$ the smallest possible common denominator, and a
primitive-necklace divisibility theorem. We isolate the exact hypothesis the
logarithmic method needs, exhibit an infinite period-three family for which
it fails modulo $5^3$, audit a printed framing statement that would have
implied the contrary, derive algebraic generating functions and two
correction terms in the asymptotic expansion, and supply exact reproducible
checks. Mathematical completeness of the written proofs and global priority
are separate matters: no claim of exhaustive novelty verification or external
peer review is made.
\end{abstract}
\vfill
\noindent\small
\textbf{Provenance.} This article merges two independently produced research
reports on OEIS A348410: \texttt{a348410-small-prime-congruences}, whose
subject was the binary completion, the denominator theorem, the M\"obius and
asymptotic material, and the source audit; and
\texttt{a348410-cubic-supercongruence}, whose subject was the odd-prime
prime-power tower, its consequences for related families, and an audit of a
broader printed framing claim. The shared odd-prime theorem
(Theorem~\ref{thm:odd}) was proved by the same argument in both reports and
is therefore proved once here. Where the two reports proved the same
statement by genuinely different routes --- the $p=3$ case of
Lemma~\ref{lem:squares}, the period-three obstruction of
Section~\ref{sec:boundary}, and the framing audit of Appendix~\ref{app:audit}
--- both arguments are preserved and marked as such.

\smallskip
\noindent\small
\textbf{Archive contents.} This article in \LaTeX\ and PDF; three verification
programs; exact CSV tables; test reports; build instructions; a source and
priority audit; BibTeX metadata; and a draft OEIS comment with a proof
outline. No computational conjecture is used as a premise of a theorem.
\end{titlepage}

\begingroup
\small
\setlength{\parskip}{0pt}
\tableofcontents
\endgroup
\clearpage

\section{The problem, the source boundary, and the results}
\label{sec:intro}

\subsection{The sequence, the conjecture, and the source boundary}
The sequence A348410 counts nonnegative integer solutions of
\begin{equation}\label{eq:baskets}
 n=\sum_{i=1}^n (u_i+2v_i).
\end{equation}
Equivalently, $n$ indistinguishable objects are distributed among $n$ labelled
white baskets and $n$ labelled black baskets, and the black occupancies are
even. We set $a(0)=1$. The OEIS entry records the coefficient formula
\begin{equation}\label{eq:an}
 a(n)=[x^n]\frac1{(1-x)^n(1-x^2)^n}
\end{equation}
and, in a comment dated 21 February 2022, Bala's conjecture
\begin{equation}\label{eq:bala}
 a(mp^r)-a(mp^{r-1})\in p^{3r}\Z
 \qquad(p\ge5\text{ prime},\ m,r\ge1).
\end{equation}
The live entry still labels this a conjecture in the version inspected for
this article \cite{oeis348410}.

That label is not a reliable certificate of present-day openness. In
particular, the 2026 articles of Niu and Prodinger already discuss the
algebraic generating function \cite{niu,prodinger}. More importantly for
this investigation, a public repository note, \emph{A cubic tower for a
two-parameter coefficient family}, contains an elementary proof candidate
for \eqref{eq:bala}, for its full integer-parameter analogue, and for a
$3^{3r-1}$ ternary bound \cite{prior}. Its parameter convention is
$[x^n]((1+x)^u(1-x)^v)^n$; ours corresponds to $(u,v)=(-\beta,-\alpha)$.
The full odd-prime parameter conjecture is also stated on A352373
\cite{oeis352373}. The repository source explicitly leaves the binary case
without an assertion.

Two questions are therefore treated here, and they have different status.
The odd-prime statement is proved in full in Section~\ref{sec:odd}, as a
first-class theorem of this article with a complete elementary proof; it is
\emph{not} claimed as a new resolution, because the prior proof candidate
exists. The binary statement is the following completion problem, which the
checked sources leave open.

\begin{problem}[Binary completion]\label{prob:binary}
For the integer-parameter family
\begin{equation}\label{eq:family}
 \A(n)=[x^n](1-x)^{-\alpha n}(1+x)^{-\beta n},
 \qquad \alpha,\beta\in\Z,
\end{equation}
does cubic growth of the divisibility exponent persist at $p=2$, with a
fixed loss independent of the level? Can one determine the first possible
binary defect and obtain an optimal common denominator for the cubic
M\"obius transform?
\end{problem}

The sources checked did not supply the binary formula or the denominator
conclusions below. This is a statement about the checked source boundary,
not an exhaustive claim about all mathematical literature. Priority for the
odd-prime result is credited to the prior proof candidate; the proof given
here is independent of it, so that the binary completion does not rest on an
external unreviewed argument.

\subsection{Why the coefficient formula is the right handle}
The displayed coefficient formula \eqref{eq:an} is far more useful for
divisibility questions than a high-order recurrence in $n$. Its logarithmic
derivative has a coefficient sequence of period two, and multiplication of
an index by an odd prime preserves parity. That single observation converts
the problem into an \emph{exact} factorization, \eqref{eq:factor} below,
rather than a congruence whose error has to be estimated. After the
factorization, the only nontrivial remaining issue is whether a weighted
reciprocal-square sum supplies one more full prime-power factor. Sections
\ref{sec:odd} and \ref{sec:boundary} show that it does in this setting, and
that the corresponding statement for period three is false.

At $p=2$ the parity argument collapses, because multiplication by $2$ does
not preserve parity. Section~\ref{sec:binary} replaces it by a reflection
identity and a formal integration by parts.

\subsection{Main theorems}
All binomial coefficients with an integer upper argument use the generalized
convention
\begin{equation}\label{eq:genbinomial}
 \binom{z}{k}=\frac{z(z-1)\cdots(z-k+1)}{k!},\qquad k\ge0.
\end{equation}
In particular, this is an integer when $z$ is a negative integer.
For an odd prime $p$ put
\begin{equation}\label{eq:epsilon}
 \eps_p=\begin{cases}1,&p=3,\\0,&p\ge5.\end{cases}
\end{equation}

\begin{theorem}[Odd-prime coefficient tower]\label{thm:odd}
Let $\alpha,\beta\in\Z$, let $p$ be an odd prime, and let $N>0$ be divisible
by $p$. Write $r=v_p(N)$. Then
\begin{equation}\label{eq:odd-tower}
 \A(N)-\A(N/p)\in p^{3r-\eps_p}\Z.
\end{equation}
\end{theorem}

Theorem~\ref{thm:odd} is proved in Section~\ref{sec:odd}. Because the
exponent is governed by the exact valuation of $N$ rather than by a declared
level, it contains the usual formulation with no coprimality restriction on
the multiplier.

\begin{corollary}[Prime-power form without coprimality]\label{cor:strong-m}
For all $m,k\ge1$ and every odd prime $p$,
\begin{equation}\label{eq:strong-m}
 \A(mp^k)\equiv\A(mp^{k-1})
 \pmod{p^{3(k+v_p(m))-\eps_p}}.
\end{equation}
In particular $a(mp^r)\equiv a(mp^{r-1})\pmod{p^{3r}}$ for $p\ge5$, which is
Bala's assertion \eqref{eq:bala}, and
$a(m3^r)\equiv a(m3^{r-1})\pmod{3^{3r-1}}$.
\end{corollary}
\begin{proof}
Apply Theorem~\ref{thm:odd} to $N=mp^k$, whose exact valuation is
$k+v_p(m)$.
\end{proof}

\begin{theorem}[First binary defect]\label{thm:binary}
Let $\alpha,\beta\in\Z$, let $m\ge1$ be odd, and let $r\ge2$. Then
\begin{align}\label{eq:binary-defect}
 &\A(m2^r)-\A(m2^{r-1})\notag\\
 &\quad\equiv
 2^{3r-3}\alpha\beta
 \binom{(\alpha+\beta+1)m-1}{m-1}
 \pmod{2^{3r-2}}.
\end{align}
Consequently the difference is always divisible by $2^{3r-3}$, and is
divisible by $2^{3r-2}$ whenever $\alpha\beta$ is even. If both parameters
are odd and $m=1$, then its valuation is exactly $3r-3$, at every $r\ge2$.
\end{theorem}

The first lift $r=1$ is deliberately not included in the displayed defect
formula. It always satisfies the elementary congruence modulo $2$ proved
below. The distinction prevents a common error: extrapolating a later-level
binary estimate to the first lift.

\begin{theorem}[All primes for the basket sequence]\label{thm:allprime}
For every prime $p$ and all $m,r\ge1$,
\begin{equation}\label{eq:all-prime}
 v_p\bigl(a(mp^r)-a(mp^{r-1})\bigr)
 \ge 3r-v_p(12).
\end{equation}
Thus the moduli are $p^{3r}$ for $p\ge5$, $3^{3r-1}$ for $p=3$, and
$2^{3r-2}$ for $p=2$. Each of these uniform losses is necessary.
\end{theorem}

For a zero difference, its valuation is interpreted as $+\infty$.

\begin{theorem}[Sharp common denominators]\label{thm:denom}
Define
\begin{equation}\label{eq:Bn}
 \B(n)=\frac1{n^3}\sum_{d\mid n}\mu(d)\A(n/d),\qquad
 B(n)=B_{2,1}(n).
\end{equation}
Then $24\B(n)\in\Z$ for all integer parameters and all $n\ge1$.
If $\alpha\beta$ is even, $12\B(n)\in\Z$. In particular,
\begin{equation}\label{eq:denom12}
 \boxed{\quad
 \frac{12}{n^3}\sum_{d\mid n}\mu(d)a(n/d)\in\Z_{\ge0}
 \quad(n\ge1).\quad}
\end{equation}
The constant $12$ is minimal for the single sequence $a(n)$. The constant
$24$ is minimal for the unrestricted integer-parameter family.
\end{theorem}

For the combinatorial interpretation, let $P_n$ be the number of primitive
cyclic words of length $n$ in the weighted alphabet of pairs $(u,v)$ of
nonnegative integers, where $(u,v)$ has weight $u+2v$, and the total word
weight is $n$. ``Primitive'' means that no shorter word is repeated to form
it. Cyclic rotations are identified.

\begin{corollary}[Primitive-orbit divisibility]\label{cor:necklace}
For every $n\ge1$,
\begin{equation}\label{eq:necklace-div}
 P_n=\frac1n\sum_{d\mid n}\mu(d)a(n/d),\qquad
 \frac{n^2}{\gcd(n^2,12)}\mid P_n.
\end{equation}
\end{corollary}

\subsection{What is proved, and what is not claimed}
All four displayed results --- Theorems~\ref{thm:odd}, \ref{thm:binary},
\ref{thm:allprime} and~\ref{thm:denom}, together with
Corollaries~\ref{cor:strong-m} and~\ref{cor:necklace} --- have complete
elementary proofs here, in Sections~\ref{sec:odd}, \ref{sec:binary}
and~\ref{sec:mobius}. They are theorems of this article, not statements
imported from elsewhere.

Novelty is a separate question, and the answer differs by result. For the
odd-prime tower, Theorem~\ref{thm:odd}, a prior public elementary proof
candidate for the same statement was located during the source audit
\cite{prior}; no priority is claimed, and the proof here is given because the
binary completion depends on it and should not rest on an external unreviewed
argument. For the binary defect formula and the sharp denominators, nothing
equivalent was found in the checked sources; that is a bounded observation
about the sources checked, not a global priority claim. The algebraic
generating function and the leading asymptotic are not claimed as new: both
are already documented in the OEIS and related literature
\cite{oeis348410,niu,prodinger}. The correction coefficients and the general
Gaussian-moment expression are derived explicitly here, without a priority
claim. That a general printed framing statement has a problem is likewise not
claimed as new; a period-four counterexample is already documented elsewhere
\cite{priorcounter}, and the two period-three witnesses of
Section~\ref{sec:boundary} and Appendix~\ref{app:audit} are compact
independent boundary evidence. A stronger exact valuation pattern at powers
of two is labelled as conjectural in Section~\ref{sec:frontiers}.

The proofs have been checked algebraically and by independent exact
computations, but have not been externally refereed or formalized in Lean.
The code is a reproducibility aid, not a substitute for the arguments. No
OEIS edit, database submission, or communication with an author has been
performed.

\section{Counting formulas and exact coefficient computation}
\label{sec:coeff}

\subsection{A positive binomial sum}
A single white basket contributes $(1-x)^{-1}$ and a black basket contributes
$(1-x^2)^{-1}$. Multiplying over the labelled baskets proves
\eqref{eq:an}. Expanding the two factors separately gives, for $n\ge1$,
\begin{equation}\label{eq:positive-sum}
 a(n)=\sum_{k=0}^{\lfloor n/2\rfloor}
 \binom{2n-2k-1}{n-2k}\binom{n+k-1}{k}.
\end{equation}
The first terms are
\[
 1,\ 1,\ 5,\ 19,\ 85,\ 376,\ 1715,\ 7890,\ 36693,\ 171820,\ 809380.
\]
Every summand in \eqref{eq:positive-sum} is nonnegative. For instance at
$n=5$ the three summands are $126$, $175$ and $75$, giving $a(5)=376$.
This positive sum gives an independent check on the faster coefficient
recurrence below. It also makes integrality and nonnegativity immediate for
the target sequence.

For arbitrary integer parameters, generalized binomial expansion shows that
\begin{equation}\label{eq:signed-sum}
 \A(n)=\sum_{k=0}^n(-1)^k
 \binom{\beta n+k-1}{k}
 \binom{\alpha n+n-k-1}{n-k}\qquad(n\ge1).
\end{equation}
This formula includes zero and negative parameters without any change of
convention. We put $\A(0)=1$.

\subsection{A two-step recurrence in the coefficient index}
Fix $n$ and write
\[
 F(x)=(1-x)^{-\alpha n}(1+x)^{-\beta n}
      =\sum_{k\ge0}f_kx^k.
\]
Logarithmic differentiation gives
\[
 (1-x^2)F'(x)=n\bigl((\alpha-\beta)+(\alpha+\beta)x\bigr)F(x).
\]
Consequently, with $f_{-1}=0$ and $f_0=1$,
\begin{equation}\label{eq:coefficient-recurrence}
 kf_k=n(\alpha-\beta)f_{k-1}
      +\bigl(n(\alpha+\beta)+k-2\bigr)f_{k-2},\qquad k\ge1.
\end{equation}
The desired output is $f_n$. This is a recurrence in $k$ \emph{with $n$
fixed}, not a recurrence relating successive diagonal values $\A(n)$.

The algorithm uses $O(n)$ exact arithmetic operations and two live integer
coefficients. The integers grow with $n$, so this is not a claim of linear
bit complexity. Every division by $k$ is checked for exactness in the
supplied implementation. Large datasets are computed without floating-point
arithmetic.

One caveat matters for anyone reusing \eqref{eq:coefficient-recurrence} in a
divisibility experiment. The division by $k$ is an exact division of
integers, and it must not be replaced by multiplication by a modular inverse
of $k$ when $p\mid k$: no such inverse exists, and silently working modulo a
prime power destroys exactly the $p$-adic information the congruences are
about.

\subsection{Small-prime boundaries visible in the first terms}
For the target sequence,
\begin{align*}
 a(2)-a(1)&=4,\\
 a(3)-a(1)&=18=2\cdot3^2,\\
 a(4)-a(2)&=80=5\cdot2^4,\\
 a(5)-a(1)&=375=3\cdot5^3,\\
 a(7)-a(1)&=7889=23\cdot7^3.
\end{align*}
These examples already show that $p=2$ and $p=3$ cannot simply be inserted
into Bala's stated modulus. They also prove the necessity of the fixed
losses in Theorem~\ref{thm:allprime}: at $p=2$ use level $r=2$, and at
$p=3$ use level $r=1$. At $p=5$ there is no extra uniform power available.
These are sharpness statements about the \emph{uniform} formulas; they do not
assert equality of the valuation at every prime, multiplier and level.

Exact valuations along the first few odd prime-power towers are collected in
Table~\ref{tab:valuations}. Every entry is computed from integers, not
inferred from a decimal approximation.
\begin{table}[htbp]
\centering
\small
\caption{Observed valuations $v_p(a(p^r)-a(p^{r-1}))$ at $m=1$.
The proved lower bound is $3r-\eps_p$.}
\label{tab:valuations}
\begin{tabular}{c r r r r}
\toprule
$p$ & $r=1$ & $r=2$ & $r=3$ & $r=4$\\
\midrule
3 & 2 & 5 & 8 & 11\\
5 & 3 & 6 & 9 & 12\\
7 & 3 & 6 & 9 & 12\\
11 & 3 & 6 & 9 & --\\
\bottomrule
\end{tabular}
\end{table}
The dash means that the index $11^4$ lies outside the default computational
sweep, not that the congruence is unknown there. As a larger explicit odd
check,
\[
 a(25)-a(5)=13\,859\,743\,127\,937\,500
\]
has $5$-adic valuation $6$, in agreement with the bound at $r=2$. The
analogous binary table is deferred to Section~\ref{sec:frontiers}, because
there the observed pattern is finer than anything proved here.

\section{Local arithmetic and the odd-prime theorem}
\label{sec:odd}

This section proves Theorem~\ref{thm:odd} in full. The argument is
elementary and self-contained: formal logarithms over $\Q$, an exact
factorization, one truncation of a formal exponential, and two coefficient
estimates. It uses neither the algebraic generating function of
Section~\ref{sec:gf} nor any general framing theorem; in particular it does
not depend on the printed rational-framing statement discussed in
Section~\ref{sec:boundary} and Appendix~\ref{app:audit}, whose explicit
conclusion is false. The same statement has a prior public proof candidate
\cite{prior}, which is credited and not superseded; the proof below is
independent of it.

\subsection{The local ring and a coefficient estimate}
For a prime $p$, let
\[
 R_p=\Z_{(p)}=\{u/v\in\Q:p\nmid v\}.
\]
Congruence modulo $p^h$ in this ring means that the difference has valuation
at least $h$. A formal power-series congruence is always coefficientwise.
All exponentials and logarithms are formal over $\Q$; at a fixed degree only
finitely many terms contribute.

Set
\begin{equation}\label{eq:phi-log}
 \Phi(x)=(1-x)^{-\alpha}(1+x)^{-\beta},\qquad
 L(x)=\log\Phi(x)=\sum_{j\ge1}\frac{c_j}{j}x^j,
 \qquad c_j=\alpha+\beta(-1)^j.
\end{equation}

\begin{lemma}[Logarithmic-derivative estimate]\label{lem:gcoef}
Let $G\in1+xR_p[[x]]$ and write
\[
 x\frac{G'(x)}{G(x)}=\sum_{i\ge1}c_ix^i,\qquad c_i\in R_p.
\]
Let $M\ge1$ and put $g_j=[x^j]G(x)^M$. Then $g_0=1$, every $g_j\in R_p$,
and for $j\ge1$,
\begin{equation}\label{eq:gvaluation}
 v_p(g_j)\ge\max\{0,v_p(M)-v_p(j)\}.
\end{equation}
If moreover $G=\Phi$ as in \eqref{eq:phi-log} with integer parameters, then
$g_j\in\Z$.
\end{lemma}
\begin{proof}
Differentiating $G^M$ and multiplying by $x$ gives
$x(G^M)'=M(xG'/G)G^M$, that is, the exact identity
\begin{equation}\label{eq:derivative-coeff}
 jg_j=M\sum_{i=1}^j c_i g_{j-i}.
\end{equation}
The sum on the right lies in $R_p$. Thus $v_p(g_j)\ge v_p(M)-v_p(j)$.
Independently $G^M\in R_p[[x]]$, which gives the lower bound zero and hence
the maximum. Integrality of $g_j$ in the case $G=\Phi$ follows from
generalized binomial expansion.
\end{proof}

The estimate is useful precisely when a complementary index is more divisible
by $p$ than $j$ itself, which is how it will be applied. No asymptotics and
no induction on a guessed recurrence are involved. Stating the lemma for a
general $G$ costs nothing and is what Remark~\ref{rem:criterion} below needs.

The second input is a truncation of the formal exponential.

\begin{lemma}[Exponential truncation]\label{lem:tail}
Let $p$ be an odd prime, let $N$ be a positive integer with
$r=v_p(N)\ge1$, and let $W\in xR_p[[x]]$. Then
\begin{equation}\label{eq:tail}
 \exp(NW)\equiv1+NW+\frac{N^2}{2}W^2
 \pmod{p^{3r-\eps_p}R_p[[x]]}.
\end{equation}
\end{lemma}
\begin{proof}
Each coefficient of $(N^k/k!)W^k$ has valuation at least $kr-v_p(k!)$.
Writing $k$ in base $p$ gives the elementary estimate
\[
 v_p(k!)=\sum_{i\ge1}\left\lfloor\frac{k}{p^i}\right\rfloor
 \le\frac{k-1}{p-1}.
\]
For $p\ge5$ this yields $v_p(k!)\le(k-1)/4\le k-3$ when $k\ge4$, and the
same bound holds directly at $k=3$ because $v_p(3!)=0$. Hence for $k\ge3$,
\[
 kr-v_p(k!)\ge3r+(k-3)(r-1)\ge3r.
\]
For $p=3$ the estimate reads $v_3(k!)\le(k-1)/2\le k-2$, so
\[
 kr-v_3(k!)\ge3r-1+(k-3)(r-1)\ge3r-1.
\]
In both cases every omitted term lies in the stated ideal, coefficient by
coefficient.
\end{proof}

\subsection{A reciprocal-square cancellation}
Recall the constant $\eps_p$ of \eqref{eq:epsilon}: it is $1$ at $p=3$ and
$0$ for $p\ge5$.

\begin{lemma}[Reduced reciprocal squares]\label{lem:squares}
If $s>0$, $p\mid s$, and $t=v_p(s)$, then
\begin{equation}\label{eq:Hs}
 H_p(s)=\sum_{\substack{1\le j<s\\p\nmid j}}\frac1{j^2}
 \in p^{t-\eps_p}R_p.
\end{equation}
\end{lemma}
\begin{proof}
Reduce denominators modulo $p^t$ and partition the summation interval into
complete blocks of length $p^t$. It is enough to consider
$U_t=\sum_{u\in(\Z/p^t\Z)^\times}u^{-2}$.

For $p\ge5$, multiplication by $2$ permutes the units and gives
$U_t\equiv 2^{-2}U_t\pmod{p^t}$. Hence $3U_t\equiv0\pmod{p^t}$; since
$3$ is invertible, $U_t\equiv0$.

For $p=3$, induction proves $U_t\equiv0\pmod{3^{t-1}}$. The case $t=1$
has modulus $1$. For $t\ge2$, each unit modulo $3^{t-1}$ has three lifts,
so $U_t\equiv3U_{t-1}\pmod{3^{t-1}}$. The induction hypothesis finishes
the proof. The same bounds hold for the full block sum.
\end{proof}

\begin{remark}[A second, uniform proof of Lemma~\ref{lem:squares}]
\label{rem:squares-alt}
The case distinction can be removed entirely. Multiplication by $2$ permutes
$(\Z/p^t\Z)^\times$ for \emph{every} odd $p$, so
\[
 U_t\equiv\sum_{a\in(\Z/p^t\Z)^\times}(2a)^{-2}
 \equiv\tfrac14U_t\pmod{p^t},
\]
whence $3U_t\equiv0\pmod{p^t}$ and therefore
$v_p(U_t)\ge t-v_p(3)=t-\eps_p$, uniformly in the odd prime $p$. What the
alternative route buys is that the exceptional prime is not an extra case at
all: the loss of one power at $p=3$ is visibly the valuation of the single
constant $3$ produced by the doubling automorphism, rather than the outcome
of a separate lifting induction. It also shows why the \emph{unweighted} sum
is safe to permute in the first place: its summand depends only on the unit
residue modulo $p^t$. The first proof is retained because its lifting
induction is the argument that generalizes when the doubling automorphism is
unavailable, as at $p=2$ in Lemma~\ref{lem:binary-H}.
\end{remark}

For odd $p$, parity is preserved by multiplication by $p$, so $c_{pj}=c_j$.
The reduced logarithm
\begin{equation}\label{eq:Dp}
 D_p(x)=L(x)-\frac1pL(x^p)
       =\sum_{\substack{j\ge1\\p\nmid j}}\frac{c_j}{j}x^j
\end{equation}
belongs to $xR_p[[x]]$ and has no monomial whose exponent is divisible by
$p$.

\begin{lemma}[Quadratic-tail estimate]\label{lem:quadratic}
If $s>0$ is divisible by the odd prime $p$, then
\begin{equation}\label{eq:Tvaluation}
 T_p(s):=[x^s]D_p(x)^2\in p^{v_p(s)-\eps_p}R_p.
\end{equation}
\end{lemma}
\begin{proof}
Write $t=v_p(s)$. Since $p\mid s$, both $j$ and $s-j$ are units whenever
$p\nmid j$. Thus
\[
 T_p(s)=\sum_{\substack{1\le j<s\\p\nmid j}}
 \frac{c_jc_{s-j}}{j(s-j)}
 \equiv-\sum_{\substack{1\le j<s\\p\nmid j}}
 \frac{c_jc_{s-j}}{j^2}\pmod{p^t}.
\]
If $s$ is odd, the weight $c_jc_{s-j}=\alpha^2-\beta^2$ is constant, so
Lemma~\ref{lem:squares} applies directly.

If $s$ is even, the weights are $(\alpha+\beta)^2$ for even $j$ and
$(\alpha-\beta)^2$ for odd $j$. The weighted reciprocal-square sum is
exactly
\[
 (\alpha-\beta)^2H_p(s)+\alpha\beta H_p(s/2).
\]
Indeed, the excess even weight is $4\alpha\beta$, while replacing $j$ by
$2k$ divides the reciprocal square by $4$. Since $p$ is odd,
$v_p(s/2)=t$. Applying Lemma~\ref{lem:squares} to both terms proves the
claim.
\end{proof}

\begin{remark}[Why the parity split, and not a permutation]
\label{rem:no-permutation}
The parity decomposition in the even case is the essential structural step,
and it cannot be replaced by an appeal to periodicity. Lemma~\ref{lem:squares}
is proved by permuting the summands of an \emph{unweighted} sum, which is
legitimate because its summand depends only on the residue of $j$ modulo
$p^t$. A weighted sum admits no such permutation: reducing an index modulo
$p^t$ can change its residue modulo the period of the weights, so a
rearrangement that fixes the reciprocal squares need not fix the weights.
What makes the step valid here is the exact identity displayed above, which
rewrites the weighted sum as a fixed linear combination of two unweighted
ones, not a rearrangement of it. Section~\ref{sec:boundary} exhibits a
period-three weight for which no such identity is available and the
conclusion genuinely fails.
\end{remark}

\subsection{Proof of the full odd-prime tower}
\begin{proof}[Proof of Theorem~\ref{thm:odd}]
Put $N=pM$, so that $v_p(M)=r-1$. From \eqref{eq:Dp} one obtains, by
splitting the logarithm at the degrees divisible by $p$, the exact identity
\begin{equation}\label{eq:factor}
 \Phi(x)^N=\Phi(x^p)^M\exp\bigl(ND_p(x)\bigr).
\end{equation}
This is an identity in $\Q[[x]]$, not a congruence whose error has to be
estimated. Lemma~\ref{lem:tail}, applied with $W=D_p$, gives
\begin{equation}\label{eq:exp-trunc}
 \exp(ND_p)\equiv1+ND_p+\frac{N^2}{2}D_p^2
 \pmod{p^{3r-\eps_p}R_p[[x]]},
\end{equation}
and thereby justifies every discarded term of the formal exponential.

Extract degree $N$. The constant term of the exponential gives $\A(M)$.
The linear contribution is zero: $\Phi(x^p)^M$ is supported on exponents
divisible by $p$, whereas $D_p$ is supported away from them. Therefore
\begin{equation}\label{eq:quadratic-reduction}
 \A(N)-\A(M)\equiv
 \frac{N^2}{2}[x^N]\Phi(x^p)^M D_p(x)^2
 \pmod{p^{3r-\eps_p}}.
\end{equation}
Writing $g_j=[x^j]\Phi(x)^M$, the remaining coefficient is
\begin{equation}\label{eq:tail-sum}
 \sum_{j=0}^{M-1}g_j T_p\bigl(p(M-j)\bigr).
\end{equation}
We claim that each summand belongs to $p^{r-\eps_p}R_p$.
For $j=0$, this follows directly from Lemma~\ref{lem:quadratic} at $s=N$.
For $j>0$ with $v_p(j)<r-1$, the equality $v_p(M)=r-1$ gives
$v_p(M-j)=v_p(j)$. Lemmas~\ref{lem:gcoef} and~\ref{lem:quadratic} then
supply, respectively,
\[
 r-1-v_p(j)\quad\hbox{and}\quad 1+v_p(j)-\eps_p,
\]
whose sum is $r-\eps_p$.
For $v_p(j)\ge r-1$, we have $v_p(M-j)\ge r-1$, and the quadratic-tail
bound alone supplies $r-\eps_p$.

Finally, $v_p(N^2/2)=2r$ because $p$ is odd. Inserting the valuation of
\eqref{eq:tail-sum} into \eqref{eq:quadratic-reduction} proves
\eqref{eq:odd-tower}. The difference is an ordinary integer, so its local
divisibility is ordinary integer divisibility.
\end{proof}

With Theorem~\ref{thm:odd} proved, Corollary~\ref{cor:strong-m} follows as
stated in Section~\ref{sec:intro}: substituting $N=mp^k$, whose exact
valuation is $k+v_p(m)$, gives \eqref{eq:strong-m} with no coprimality
restriction on $m$.

\subsection{Where the three powers come from}
\label{sec:three-powers}
The exponent $3r$ is not a numerical coincidence, and it is worth recording
how it is assembled. Two of the three powers of $p$ come from the scalar
$N^2/2$ in front of the quadratic exponential term. The third is produced
jointly by the two factors of a single summand of \eqref{eq:tail-sum}. For
the summand indexed by $j$ with $t=v_p(j)<r-1$, the quadratic-tail estimate
of Lemma~\ref{lem:quadratic} at $s=p(M-j)$ contributes $1+t-\eps_p$ powers,
which is less than $r-\eps_p$; exactly the shortfall $r-1-t$ is then
supplied by the coefficient $g_j$ through Lemma~\ref{lem:gcoef}. The two
estimates are complementary: whichever of the two indices $j$ and $M-j$
carries little $p$-divisibility, the other carries the rest.

This is also the reason a proof written only for $N=p$ would miss an
essential part of the argument. At $N=p$ one has $M=1$, the sum
\eqref{eq:tail-sum} has the single term $j=0$, and the coefficient estimate
of Lemma~\ref{lem:gcoef} is never used. The interaction between the two
estimates, which is what makes the tower climb, appears only at level two
and beyond.

\begin{remark}[A sufficient criterion beyond this family]
\label{rem:criterion}
Fix a prime $p\ge5$. The proof above used two properties of $\Phi$ and
nothing else. Accordingly, let $G\in1+xR_p[[x]]$ be any series whose
logarithmic-derivative coefficients $xG'/G=\sum_{j\ge1}c_jx^j$ satisfy
$c_{pj}=c_j$, and put $A(n)=[x^n]G(x)^n$. If in addition
\begin{equation}\label{eq:criterion}
 [x^s]\Bigl(\sum_{p\nmid j}\frac{c_j}{j}x^j\Bigr)^2
 \in p^{v_p(s)}R_p
 \qquad\text{for every }s\text{ with }p\mid s,
\end{equation}
then $A(N)-A(N/p)\in p^{3v_p(N)}R_p$ for every $N$ divisible by $p$.
The proof is verbatim that of Theorem~\ref{thm:odd}, with
Lemma~\ref{lem:quadratic} replaced by the hypothesis \eqref{eq:criterion}
and Lemma~\ref{lem:gcoef} applied to $G$ in its general form.

The point of isolating the criterion is the word ``in addition''. Exact
invariance $c_{pj}=c_j$ under index multiplication gives the factorization
and kills the linear term, but it does not give \eqref{eq:criterion}:
the quadratic estimate is a genuinely additional hypothesis, not a
consequence of periodicity or of Frobenius invariance.
Section~\ref{sec:boundary} exhibits a period-three sequence satisfying
$c_{5j}=c_j$ for which \eqref{eq:criterion} fails at $p=5$, and with it the
conclusion.
\end{remark}

\subsection{Other integer-sequence families}
\label{sec:families}
Theorem~\ref{thm:odd} covers more than the basket sequence, and a few
reparametrizations make the range visible.

Since $(1-x)(1+x)=1-x^2$, a change of parameters turns the family into a
two-kind basket family. For integers $\sigma,\tau$,
\begin{equation}\label{eq:alpha-beta}
 A_{\sigma+\tau,\,\tau}(n)
 =[x^n](1-x)^{-\sigma n}(1-x^2)^{-\tau n}.
\end{equation}
For nonnegative $\sigma,\tau$ the right side counts distributions of $n$
objects into $\sigma n$ labelled unrestricted baskets and $\tau n$ labelled
even-occupancy baskets; the case $\sigma=\tau=1$ is $A_{2,1}=a$, that is,
\eqref{eq:baskets}. All of these sequences obey the same odd-prime
congruences. Negative $\sigma$ or $\tau$ remains covered algebraically
through \eqref{eq:signed-sum}, but the basket reading is then unavailable.

Two elementary specializations are useful as sanity checks. For a positive
integer $\alpha$,
\begin{equation}\label{eq:spec-u0}
 A_{\alpha,0}(n)=\binom{(\alpha+1)n-1}{n}\qquad(n\ge1),
\end{equation}
by the generalized hockey-stick identity. For a positive integer $q$,
$\Phi=(1-x^2)^{-q}$ when $\alpha=\beta=q$, so
\begin{equation}\label{eq:spec-qq}
 A_{q,q}(n)=
 \begin{cases}
 0,&n\text{ odd},\\[2pt]
 \displaystyle\binom{qn+n/2-1}{n/2},&n\text{ even}.
 \end{cases}
\end{equation}
Theorem~\ref{thm:odd} accommodates the vanishing odd terms and the nonzero
even ones without a separate divisibility convention, since $0$ is divisible
by every power of $p$.

Finally, the step of the second factor may be changed freely away from its
own prime divisors.

\begin{corollary}[Change of step]\label{cor:step}
Let $d\ge1$ and $\alpha,\beta\in\Z$, and put
\[
 A^{(d)}(n)=[x^n](1-x^d)^{-\alpha n}(1+x^d)^{-\beta n}.
\]
Let $p$ be an odd prime with $p\nmid d$, and let $N>0$ be divisible by $p$.
Then
\[
 A^{(d)}(N)-A^{(d)}(N/p)\in p^{3v_p(N)-\eps_p}\Z.
\]
\end{corollary}
\begin{proof}
The series $(1-x^d)^{-\alpha n}(1+x^d)^{-\beta n}$ is supported on exponents
divisible by $d$. Since $p\nmid d$, the conditions $d\mid N$ and
$d\mid N/p$ are equivalent. If they fail, both coefficients vanish and the
difference is $0$. If they hold, then
$A^{(d)}(N)=A_{\alpha,\beta}(N/d)$ and $A^{(d)}(N/p)=A_{\alpha,\beta}(N/(pd))$,
while $v_p(N/d)=v_p(N)$ because $p\nmid d$. Theorem~\ref{thm:odd} applied at
$N/d$ gives the claim.
\end{proof}

The hypothesis $p\nmid d$ is not decorative, and the restriction to the
divisors of $d$ is not the only one that matters: Section~\ref{sec:boundary}
shows that replacing $1+x^2$ by $1-x^3$ --- a change of \emph{shape} rather
than of step --- destroys the conclusion at $p=5$, a prime that divides
neither the old period nor the new one.

\subsection{Limits along a prime-power tower}
\label{sec:padic-limit}
Fix an odd prime $p$, integer parameters, and a positive integer $m$. By
Corollary~\ref{cor:strong-m} the differences
$\A(mp^{r})-\A(mp^{r-1})$ have valuations tending to infinity, so the
sequence $r\mapsto\A(mp^r)$ is Cauchy in $\Z_p$ and converges to a limit
$\lambda_{\alpha,\beta,p,m}\in\Z_p$. Summing the tail of successive
differences gives the error estimate
\begin{equation}\label{eq:padic-limit}
 v_p\bigl(\lambda_{\alpha,\beta,p,m}-\A(mp^r)\bigr)
 \ge3\bigl(r+1+v_p(m)\bigr)-\eps_p\qquad(r\ge0).
\end{equation}
This is a consequence of the full tower, not of the first congruence alone,
and it does not require identifying the limit in closed form. We make no
claim here about what $\lambda_{\alpha,\beta,p,m}$ is.

\section{The binary defect: proof and sharpness}
\label{sec:binary}

At $p=2$, the equality $c_{pj}=c_j$ used in \eqref{eq:Dp} fails. The correct
replacement is a reflection argument. It separates the even and odd parts
before any exponential is truncated.

\subsection{An exact odd-harmonic normalization}
Work in $R_2=\Z_{(2)}$ and define
\begin{equation}\label{eq:U}
 U(x)=\sum_{\substack{j\ge1\\j\text{ odd}}}\frac{x^j}{j}
     =\frac12\log\frac{1+x}{1-x}.
\end{equation}
The factor $1/2$ in the logarithmic expression is harmless: the displayed
odd-power series has coefficients in $R_2$.

\begin{lemma}[Odd harmonic sums are exactly normalized]\label{lem:binary-H}
For every positive even $T$, let
\[
 h(T)=\sum_{\substack{1\le j<T\\j\text{ odd}}}\frac1j.
\]
Then
\begin{equation}\label{eq:h-exact}
 v_2(h(T))=2v_2(T)-2.
\end{equation}
In particular, $h(2j)/j^2$ is an odd element of $R_2$ for every $j\ge1$.
\end{lemma}
\begin{proof}
Write $T=2^e q$ with $q$ odd. First,
\begin{equation}\label{eq:unit2}
 \sum_{\substack{1\le u<2^e\\u\text{ odd}}}u^{-2}
 \equiv2^{e-1}\pmod{2^e}.
\end{equation}
For $e=1$ this is immediate. For $e\ge2$, the two lifts $u$ and
$u+2^{e-1}$ of an odd residue modulo $2^{e-1}$ have equal inverse squares
modulo $2^e$. Indeed their squares differ by a multiple of $2^e$.
Thus the sum at level $e$ is twice the sum at level $e-1$ modulo $2^e$,
which proves \eqref{eq:unit2} by induction.

Partitioning the interval $1\le j<T$ into $q$ blocks now shows that
$\sum_{j\text{ odd}}j^{-2}\equiv2^{e-1}\pmod{2^e}$. Pair $j$ with
$T-j$ to obtain the exact identity
\[
 h(T)=\frac T2
 \sum_{\substack{1\le j<T\\j\text{ odd}}}\frac1{j(T-j)}.
\]
Modulo $2^e$, the sum on the right is the negative of the reciprocal-square
sum. It consequently has valuation exactly $e-1$. The prefactor $T/2$
has valuation $e-1$ as well, proving \eqref{eq:h-exact}.
\end{proof}

The even power series $U(x)^2$ satisfies
\[
 [x^{2j}]U(x)^2
 =\sum_{\substack{1\le k<2j\\k\text{ odd}}}\frac1{k(2j-k)}
 =\frac{h(2j)}j.
\]
Define
\begin{equation}\label{eq:K}
 K(t)=\sum_{j\ge1}\frac{h(2j)}{j^2}t^j\in tR_2[[t]],
 \qquad V(t)=U(\sqrt t)^2.
\end{equation}
Here $U(\sqrt t)^2$ denotes the integer-exponent series obtained from the
even series $U(x)^2$ by replacing $x^2$ with $t$; no choice of an analytic
square root is involved. The name $V$ is local to this section and is
unrelated to the framing inputs $V_1$ and $V_3$ of
Section~\ref{sec:boundary} and Appendix~\ref{app:audit}. The identities we need are
\begin{equation}\label{eq:K-identities}
 V(t)=tK'(t),\qquad
 U(t)^2=\frac t2\frac{\mathrm d}{\mathrm dt}K(t^2),\qquad
 K(t)\equiv\frac{t}{1-t}\pmod2.
\end{equation}
The first two follow by coefficient comparison. The last follows from the
oddness assertion of Lemma~\ref{lem:binary-H}. This final residue, not just
a divisibility bound, will determine the leading binary digit.

\subsection{Separating the two coefficient problems}
Set
\[
 S=\alpha+\beta,\qquad \delta=\alpha-\beta.
\]
Let $M$ be even, put $s=v_2(M)\ge1$, and define
\[
 g(t)=(1-t)^{-SM}.
\]
The logarithmic identities
\begin{align*}
 \Phi(x)^{2M}&=(1-x^2)^{-SM}\exp\bigl(2\delta M U(x)\bigr),\\
 \Phi(t)^M&=g(t)\exp\bigl(-2\beta M U(t)\bigr)
\end{align*}
are exact. Only the even part of the first exponential contributes to degree
$2M$. Consequently
\begin{align}\label{eq:binary-factor}
 \A(2M)&=[t^M]g(t)\cosh\bigl(2\delta M U(\sqrt t)\bigr),\notag\\
 \A(M)&=[t^M]g(t)\exp\bigl(-2\beta M U(t)\bigr).
\end{align}
The hyperbolic cosine is interpreted as its even formal Taylor series,
so every exponent of $t$ in \eqref{eq:binary-factor} is an integer.

Since $v_2(k!)\le k-1$, for $k\ge3$ we have
\[
 v_2\left(\frac{(2M)^k}{k!}\right)\ge ks+1\ge3s+1.
\]
The same estimate applies to all even degrees $k\ge4$ in the cosine.
Thus modulo $2^{3s+1}$ we need only three coefficient expressions:
\begin{equation}\label{eq:I-def}
 I_1=[t^M]gU,\qquad I_2=[t^M]gU^2,\qquad J=[t^M]gV.
\end{equation}
Writing $\Delta=\A(2M)-\A(M)$, expansion of
\eqref{eq:binary-factor} gives
\begin{equation}\label{eq:delta-start}
 \Delta\equiv
 2\beta M I_1+2\delta^2M^2J-2\beta^2M^2I_2
 \pmod{2^{3s+1}}.
\end{equation}

\subsection{Reflection converts the linear term into a quadratic one}
The apparent obstacle in \eqref{eq:delta-start} is the linear term $I_1$.
Because $M$ is even and $U$ is odd,
\[
 [t^M]g(-t)U(t)=-I_1.
\]
Also
\[
 \frac{g(-t)}{g(t)}=
 \left(\frac{1-t}{1+t}\right)^{SM}
 =\exp\bigl(-2SMU(t)\bigr).
\]
Combining these identities yields
\begin{equation}\label{eq:reflection}
 I_1=\frac12[t^M]g(t)U(t)
       \left(1-\exp\bigl(-2SMU(t)\bigr)\right).
\end{equation}
The first exponential term in \eqref{eq:reflection} is $SMI_2$.
Every remaining term has a scalar coefficient with valuation at least
$2s$: for $k\ge2$,
\[
 v_2\left(\frac{(2SM)^k}{2k!}\right)
 \ge k(s+1)-1-v_2(k!)\ge ks\ge2s.
\]
Hence
\begin{equation}\label{eq:I1-reduction}
 I_1\equiv SMI_2\pmod{2^{2s}R_2}.
\end{equation}
Multiplication by $2\beta M$ makes the error divisible by $2^{3s+1}$.
Substituting in \eqref{eq:delta-start} and using $\beta S-\beta^2=\alpha\beta$
therefore gives
\begin{equation}\label{eq:delta-quadratic}
 \Delta\equiv2M^2\bigl(\alpha\beta I_2+\delta^2J\bigr)
 \pmod{2^{3s+1}}.
\end{equation}
This reflection is the step that replaces the odd-prime support cancellation.
It is valid because $M$ is even; this is exactly why the first binary lift
was separated from the theorem.

\subsection{Formal integration by parts supplies the third power}
For any $K_0(t)\in R_2[[t]]$, coefficient differentiation gives
\begin{equation}\label{eq:ibp}
 [t^M]g(t)tK_0'(t)
 =M[t^M]g(t)K_0(t)\left(1-\frac{St}{1-t}\right).
\end{equation}
Indeed, $[t^M]t(gK_0)'=M[t^M]gK_0$ and
$tg'/g=SMt/(1-t)$.

Using \eqref{eq:K-identities} in \eqref{eq:ibp} shows
\begin{equation}\label{eq:IJ-div}
 J\in MR_2,\qquad
 I_2=\frac M2 C,
\end{equation}
where
\begin{equation}\label{eq:C}
 C=[t^M](1-t)^{-SM}K(t^2)
       \left(1-\frac{St}{1-t}\right)\in R_2.
\end{equation}
Thus the $J$ term in \eqref{eq:delta-quadratic} is already zero modulo
$2^{3s+1}$, and we obtain the particularly useful congruence
\begin{equation}\label{eq:delta-C}
 \Delta\equiv\alpha\beta M^3 C\pmod{2^{3s+1}}.
\end{equation}
It immediately proves the universal valuation bound $3s$ and its one-bit
improvement when $\alpha\beta$ is even.

It remains to compute $C$ modulo $2$ when both parameters are odd.
In this case $S$ is even. Write $M=2^s m$ with $m$ odd.
Over $\mathbb F_2[[t]]$, generalized integer powers and the Frobenius identity
give
\[
 (1-t)^{-SM}\equiv(1-t^{2^s})^{-Sm},\qquad
 1-\frac{St}{1-t}\equiv1,
\]
while $K(t^2)\equiv t^2/(1-t^2)$. Since $2^s\ge2$, extraction of $t^M$
then gives
\begin{align}\label{eq:C-binomial}
 C&\equiv\sum_{j=0}^{m-1}\binom{Sm+j-1}{j}\notag\\
  & =\binom{(S+1)m-1}{m-1}\pmod2.
\end{align}
The last identity is the generalized hockey-stick identity, obtained by
multiplying $(1-z)^{-Sm}$ by $(1-z)^{-1}$ and extracting degree $m-1$.
It holds for every integer $S$, including negative values.

If $\alpha\beta$ is even, both sides of the desired leading-bit formula
vanish regardless of $C$. If $\alpha\beta$ is odd, insert
\eqref{eq:C-binomial} into \eqref{eq:delta-C}. Finally,
$M^3=2^{3s}m^3\equiv2^{3s}\pmod{2^{3s+1}}$.
Taking $s=r-1$ proves Theorem~\ref{thm:binary} in full. For $m=1$ the binomial
coefficient in \eqref{eq:binary-defect} equals $1$, so odd parameters give
an odd leading quotient and exact valuation $3r-3$.

\subsection{First lift, optimal losses, and the all-prime statement}
For any $\Phi\in1+x\Z[[x]]$,
\[
 \Phi(x)^2\equiv\Phi(x^2)\pmod2.
\]
Raising this to $m$ and extracting degree $2m$ proves
\begin{equation}\label{eq:firstlift}
 \A(2m)\equiv\A(m)\pmod2.
\end{equation}
For $(\alpha,\beta)=(2,1)$, the improved binary bound of
Theorem~\ref{thm:binary} is $3r-2$ for $r\ge2$, and
\eqref{eq:firstlift} supplies exactly the claimed lower bound $1$ at $r=1$.
If $m$ is even, apply the theorem at the larger exact binary level.
Together with Theorem~\ref{thm:odd}, this proves
Theorem~\ref{thm:allprime}. The examples in Section~\ref{sec:coeff} prove
its uniform sharpness.

There is a distinction between uniform sharpness and eventual sharpness.
For the target sequence, equality occurs at $r=2$, but later powers of two
show more divisibility. The theorem does not assert equality at every level.
For the unrestricted family, in contrast, odd parameters and $m=1$ attain
the general bound $3r-3$ at every level $r\ge2$.

\section{M\"obius transforms, common denominators, and cyclic words}
\label{sec:mobius}

\subsection{From a prime-power tower to a denominator theorem}
For any sequence $A(n)$, set
\[
 T_A(n)=\sum_{d\mid n}\mu(d)A(n/d).
\]
If $n=p^r n_0$ with $p\nmid n_0$, separating divisors according as they
contain $p$ gives the exact identity
\begin{equation}\label{eq:mobius-group}
 T_A(n)=\sum_{d\mid n_0}\mu(d)
 \left(A(p^r n_0/d)-A(p^{r-1}n_0/d)\right).
\end{equation}
Only squarefree divisors matter, so there are no terms containing a higher
power of $p$ in the divisor itself.

Apply \eqref{eq:mobius-group} to the basket sequence. Theorem~\ref{thm:allprime}
implies, for every prime $p\mid n$,
\[
 v_p(T_a(n))\ge3v_p(n)-v_p(12).
\]
Therefore $n^3\mid12T_a(n)$, proving the integrality part of
\eqref{eq:denom12}.

For the general family, the odd-prime bounds are unchanged and the binary
bound loses at most three powers. At level one the needed general estimate
is merely valuation at least zero. Consequently
\[
 n^3\mid24T_{\A}(n).
\]
If $\alpha\beta$ is even, the improved binary estimate and
\eqref{eq:firstlift} reduce $24$ to $12$. This proves all the integrality
assertions of Theorem~\ref{thm:denom}.

The constants cannot be reduced. For the target sequence,
\begin{equation}\label{eq:denom-sharp}
 B(3)=\frac{19-1}{27}=\frac23,
 \qquad B(4)=\frac{85-5}{64}=\frac54.
\end{equation}
Every integer clearing all denominators must thus be divisible by both $3$
and $4$. Their least common multiple is $12$. For the unrestricted family,
$(\alpha,\beta)=(1,1)$ gives
\[
 B_{1,1}(4)=\frac{10-2}{64}=\frac18,
\]
and the target example at $n=3$ forces a factor $3$, so a universal clearing
integer must be divisible by $24$. Even the single parameter pair $(1,3)$
requires $24$: its values at $n=3$ and $n=4$ are $-8/3$ and $59/8$.

\subsection{Primitive cyclic words and positivity}
Let $\mathcal A=\{(u,v):u,v\in\Z_{\ge0}\}$, with weight $u+2v$.
A word of length $n$ and total weight $n$ corresponds exactly to a solution
of \eqref{eq:baskets}. Although the alphabet is infinite, the weight
constraint permits only finitely many letters at each fixed $n$.

A word of length $n$ whose least period is $d\mid n$ is the repetition of
a primitive block of length $d$. Its total weight is $n$ precisely when the
block weight is $d$. A primitive cyclic orbit of such blocks has exactly
$d$ distinct linear representatives. Therefore
\begin{equation}\label{eq:period}
 a(n)=\sum_{d\mid n}dP_d.
\end{equation}
M\"obius inversion proves the first formula of \eqref{eq:necklace-div} and
also proves $P_n\ge0$. Since $B(n)=P_n/n^2$, nonnegativity in
Theorem~\ref{thm:denom} follows. Integrality of $12P_n/n^2$ is equivalent
to the second assertion of \eqref{eq:necklace-div}.

The first values illustrate that this is stronger than the ordinary
integrality of a necklace count:
\begin{center}
\begin{tabular}{r r r r}
\toprule
$n$ & $P_n$ & $B(n)=P_n/n^2$ & $12B(n)$\\
\midrule
1&1&$1$&12\\
2&2&$1/2$&6\\
3&6&$2/3$&8\\
4&20&$5/4$&15\\
5&75&$3$&36\\
6&282&$47/6$&94\\
7&1127&$23$&276\\
8&4576&$143/2$&858\\
9&19089&$707/3$&2828\\
10&80900&$809$&9708\\
11&348238&$2878$&34536\\
12&1516620&$126385/12$&126385\\
\bottomrule
\end{tabular}
\end{center}
For example, $P_7=49\cdot23$, and $P_6$ is divisible by
$36/\gcd(36,12)=3$. The individual counts need not be multiples of $n^2$;
small-prime denominators are genuine.

The entry at $n=12$ deserves a remark. The minimality argument above uses
two indices, $n=3$ and $n=4$, to force the factors $3$ and $4$ separately.
The value $B(12)=126385/12$ shows more: the full denominator $12$ is already
attained at a single index, so no proper divisor of $12$ clears the
transform even along one arithmetic progression.

\subsection{Lambert series and an Euler product}
Let $H(z)=\sum_{n\ge0}a(n)z^n$. Equation~\eqref{eq:period} gives the formal
Lambert expansion
\begin{equation}\label{eq:Lambert}
 H(z)-1=\sum_{n\ge1}nP_n\frac{z^n}{1-z^n}
       =\sum_{n\ge1}n^3B(n)\frac{z^n}{1-z^n}.
\end{equation}
All operations are valid formally because only finitely many divisors
contribute to a fixed coefficient. The denominator theorem can be stated
as follows: the cubic Lambert coefficients of $12(H-1)$ are nonnegative
integers, and $12$ is the least positive integer with this property.

A related Euler product will follow from the Lagrange series in the next
section:
\begin{equation}\label{eq:euler-product}
 \frac{y(z)}z=\prod_{n\ge1}(1-z^n)^{-P_n}.
\end{equation}
To verify it directly, take formal logarithms. The coefficient of $z^k$
on the right is $k^{-1}\sum_{n\mid k}nP_n=a(k)/k$, exactly the logarithm
on the left derived below.

\section{Algebraic generating functions}
\label{sec:gf}

The main generating-function identities are established facts for A348410.
Neither the parametric generating function, nor the quartic it satisfies,
nor the extraction recurrence \eqref{eq:coefficient-recurrence} is claimed as
new: the first two are recorded in the OEIS entry \cite{oeis348410} and
treated by Niu \cite{niu} and Prodinger \cite{prodinger}, and Prodinger also
extracts a recurrence. We include short derivations because they connect the
coefficient-power definition, the Euler product, and the asymptotic analysis,
and because they make the symbolic checks transparent.

It is worth saying plainly what this section is \emph{not} for.
Algebraicity is not used to prove any supercongruence in this article. The
logarithmic coefficient pattern of $\Phi$, rather than the algebraicity of
$H$, is the input that supplies the divisibility; Remark~\ref{rem:criterion}
and Section~\ref{sec:boundary} show that the relevant hypothesis is a
statement about $\log\Phi$, and that algebraic sequences of the same shape
can fail it.

\subsection{Lagrange reversion for the full parameter family}
Let $y(z)$ be the unique series in $z\Z[[z]]$ satisfying
\begin{equation}\label{eq:reversion}
 y=z\Phi(y).
\end{equation}
Existence and uniqueness follow recursively from $\Phi(0)=1$. Write
$H_{\alpha,\beta}(z)=\sum_{n\ge0}\A(n)z^n$.
Lagrange inversion, applied to $\log\Phi(y)$, gives
\begin{align*}
 [z^n]\log\Phi(y(z))
 &=\frac1n[x^{n-1}]L'(x)\Phi(x)^n\\
 &=\frac1{n^2}[x^{n-1}]\frac{\mathrm d}{\mathrm dx}\Phi(x)^n
 =\frac{\A(n)}n.
\end{align*}
Since $y/z=\Phi(y)$, it follows that
\begin{equation}\label{eq:log-y}
 \log\frac{y(z)}z=\sum_{n\ge1}\frac{\A(n)}n z^n.
\end{equation}
Differentiating and using \eqref{eq:reversion} yields
\begin{equation}\label{eq:Hparam}
 H_{\alpha,\beta}(z)=\frac{zy'(z)}{y(z)}
 =\frac{1-y^2}
 {1-(\alpha-\beta)y-(\alpha+\beta+1)y^2}.
\end{equation}
For integer parameters, $\Phi$ is rational. Thus \eqref{eq:reversion}
becomes a polynomial equation after clearing denominators, and
$H_{\alpha,\beta}$ is algebraic.

\subsection{The target quartic}
For $(\alpha,\beta)=(2,1)$,
\begin{equation}\label{eq:target-param}
 z=y(1-y)^2(1+y)=y-y^2-y^3+y^4,\qquad
 H(z)=\frac{1-y^2}{1-y-4y^2}.
\end{equation}
The expanded form on the left has linear coefficient $1$, so the formal
inverse $y(z)\in z\Z[[z]]$ exists and is unique; this is the uniqueness
already noted after \eqref{eq:reversion}, made explicit for the target.
Eliminating $y$ gives
\begin{equation}\label{eq:quartic}
 \begin{split}
 z^2(4H-1)^4
 +zH(107H^3-107H^2+36H-4)\\
 +32H^3(1-H)=0.
 \end{split}
\end{equation}
Substitution of \eqref{eq:target-param} verifies the identity directly.
The partial derivative of the left side with respect to $H$, at
$(z,H)=(0,1)$, is $-32$. Thus the formal implicit-function theorem singles
out the unique branch with $H(0)=1$. In particular, no ambiguous choice of
an algebraic root is hidden in the generating-function description.

Collecting \eqref{eq:quartic} by powers of $H$ instead of by powers of $z$
gives the equivalent grouped form
\begin{align}\label{eq:quartic-grouped}
 0={}&(256z^2+107z-32)(H^4-H^3)\notag\\
    &+(96z^2+36z)H^2-(16z^2+4z)H+z^2,
\end{align}
which is the shape in which the elimination is easiest to check by hand.
Expanding both displays shows them to be the same polynomial; the symbolic
program verifies this as well.

The symbolic verification programs check the elimination in two independent
ways --- that the resultant of the two defining equations equals the
displayed quartic, and that the quotient of the resultant by it is a nonzero
constant, recorded as $1$ --- together with the parametrized substitution and
the series identity through degree $20$. The exact-arithmetic program
independently verifies, in integer arithmetic alone and to degree $60$, the
inverse defining polynomial $y-y^2-y^3+y^4=z$, the parametric form
$(1-y-4y^2)H=1-y^2$, and the quartic itself. The coefficient recurrence
\eqref{eq:coefficient-recurrence} supplies the series in both cases, so
neither check merely substitutes one rearranged symbolic expression into
another.

\section{A boundary test for periodic-coefficient arguments}
\label{sec:boundary}

\subsection{What the logarithmic method actually requires}
The odd-prime proof used two different ingredients: the equality $c_{pj}=c_j$
and the much stronger quadratic cancellation in Lemma~\ref{lem:quadratic}.
The first does not imply the second. Remark~\ref{rem:criterion} states the
exact sufficient hypothesis; this section shows that it is not automatic.

In detail, let $G\in1+xR_p[[x]]$ and suppose
\[
 \log G(x)=\sum_{j\ge1}c_jx^j/j,\qquad c_j\in R_p,\qquad c_{pj}=c_j.
\]
Define $A(n)=[x^n]G(x)^n$ and the reduced logarithm as in \eqref{eq:Dp}. The
factorization and the vanishing linear term still hold. For odd $p$, all
terms of degree at least two in the exponential have valuation at least
$2v_p(N)$, so they yield a quadratic tower. A cubic tower requires the
additional estimate \eqref{eq:criterion} on the coefficients of the square
of the reduced logarithm.

At the first layer and for $p\ge5$, this gives a particularly transparent
necessary and sufficient residue condition:
\begin{equation}\label{eq:obstruction}
 A(p)-A(1)\equiv
 -\frac{p^2}{2}\sum_{j=1}^{p-1}\frac{c_jc_{p-j}}{j^2}
 \pmod{p^3}.
\end{equation}
Indeed, the quadratic coefficient at degree $p$ is unaffected by positive
powers of $x^p$ from $\Phi(x^p)$, and
$1/(j(p-j))\equiv-1/j^2\pmod p$. This proves the formula without any
periodicity assumption beyond $c_{pj}=c_j$.

\subsection{A spacing-three counterexample}
Replace the even-occupancy restriction by occupancy divisible by $3$. More
generally, for a positive integer $t$ put
\begin{equation}\label{eq:bn-boundary}
 b_t(n)=[x^n](1-x)^{-tn}(1-x^3)^{-tn},\qquad b(n)=b_1(n).
\end{equation}
The logarithmic coefficients of $(1-x)^{-t}(1-x^3)^{-t}$ are
$c_j=t(1+3\mathbf1_{3\mid j})$. They have period three and satisfy
$c_{5j}=c_j$, so the exact logarithmic splitting of Section~\ref{sec:odd}
still works at $p=5$ and the linear term still vanishes. The failure is
elsewhere. Nevertheless, at $t=1$,
\begin{align*}
 b(1)&=1,\\
 b(5)&=\binom95+5\binom62=126+75=201,\\
 b(5)-b(1)&=200\not\equiv0\pmod{125}.
\end{align*}
Note that $5\nmid3$: the failure is not caused by the prime dividing the
period.

\paragraph{First witness: the obstruction formula.}
The obstruction in \eqref{eq:obstruction} is exactly
\begin{equation}\label{eq:weighted-counterexample}
 \sum_{j=1}^4\frac{c_jc_{5-j}}{j^2}
 =1+\frac44+\frac49+\frac1{16}
 =\frac{361}{144}\equiv4\pmod5.
\end{equation}
Equation~\eqref{eq:obstruction} predicts the residue $75$ modulo $125$,
which agrees with $200\bmod125$.

\paragraph{Second witness: the quadratic coefficient itself.}
\label{par:direct}
The same failure can be certified without passing through
\eqref{eq:obstruction}, by computing the quantity that
Remark~\ref{rem:criterion} actually requires. At $t=1$ and $p=5$,
\begin{equation}\label{eq:direct-quadratic}
 [x^5]\Bigl(\sum_{5\nmid j}\frac{c_j}{j}x^j\Bigr)^2
 =2\left(1\cdot\frac14+\frac12\cdot\frac43\right)
 =\frac{11}{6}\not\equiv0\pmod5,
\end{equation}
the two products coming from the index pairs $(1,4)$ and $(2,3)$. Since
$v_5(5)=1$, hypothesis \eqref{eq:criterion} demands divisibility by $5$ and
is violated. What this second route buys is directness: it exhibits the
failure in the precise quantity the proof consumes, with no appeal to the
first-level residue formula \eqref{eq:obstruction} and hence no reliance on
the reduction $1/(j(p-j))\equiv-1/j^2$. The two computations are consistent:
the quantity reduced in \eqref{eq:Tvaluation} is the negative of the one in
\eqref{eq:weighted-counterexample}, and indeed $11/6$ and $-361/144$ are both
congruent to $1$ modulo $5$.

\paragraph{An infinite family of first-level failures.}
The example is not isolated. For $t$ with $5\nmid t$, the factor
$(1-x^3)^{-5t}$ can contribute only its constant term or its term
$5tx^3$ in degree $5$, so
\[
 b_t(5)=\binom{5t+4}{5}+5t\binom{5t+1}{2},\qquad b_t(1)=t.
\]
Theorem~\ref{thm:odd} applied at $(\alpha,\beta)=(t,0)$, where
\eqref{eq:spec-u0} gives $A_{t,0}(5)=\binom{5t+4}{5}$ and $A_{t,0}(1)=t$,
yields $\binom{5t+4}{5}\equiv t\pmod{125}$. Therefore
\begin{equation}\label{eq:period3-family}
 b_t(5)-b_t(1)\equiv5t\binom{5t+1}{2}=\frac{25t^2(5t+1)}{2}
 \equiv\frac{25t^2}{2}\pmod{125},
\end{equation}
the denominator $2$ being read as a unit modulo $125$. Hence
$v_5(b_t(5)-b_t(1))=2$ exactly, for every $t$ not divisible by $5$. The
members $t=1,3,9$ give the differences $200$, $13\,425$ and $1\,953\,450$,
with residues $75$, $50$ and $75$ modulo $125$. The member $t=3$ is the one
needed in Appendix~\ref{app:audit}, because its logarithmic coefficient
sequence satisfies the rational $2$-function conditions at \emph{every}
prime, including $3$.

Thus the success of period two is structural, not an automatic consequence
of rationality, periodicity, or Frobenius invariance. In the period-two
proof, the weighted square sum collapses to two ordinary reciprocal-square
sums, as Remark~\ref{rem:no-permutation} explains. That special collapse is
absent here.

\subsection{Scope of the literature correction}
\label{sec:litscope}
M\"uller's arXiv version 1 of \emph{Wolstenholme Type Congruences and Framing
of Rational 2-Functions} states a general periodic weighted-harmonic bound
in Theorems 1.2 and 6.2, a corresponding precise cubic framing bound in
Theorem 1.1, and the framing coefficient formula in Proposition 5.2(3)
\cite{muller}. The value \eqref{eq:weighted-counterexample}
contradicts the printed harmonic bound at period $3$ and prime $5$.

The coefficient counterexample also fits the printed number-field
hypotheses, not merely a superficial formal analogy. Take
$K=\Q(\sqrt{-3})$, whose discriminant is $-3$, and
\begin{equation}\label{eq:V1-input}
 V_1(x)=\frac{x}{1-x}+\frac{3x^3}{1-x^3},
 \qquad c_j=1+3\mathbf1_{3\mid j}.
\end{equation}
At every unramified rational prime $p\ne3$, its rational coefficients
satisfy $c_{mp^r}=c_{mp^{r-1}}$ exactly. Thus it is a rational 2-function
in the convention of that preprint. Proposition 5.2(3) there identifies the
coefficient of its plus-framing with parameter one as
$[x^n]\exp(n\int V_1)=b(n)$, where the formal integration sends $x^j$ to
$x^j/j$. Prime $5$ is unramified and does not divide the period. The
precise bound in Theorem 1.1 would therefore give divisibility by $125$,
contrary to $b(5)-b(1)=200$.

Appendix~\ref{app:audit} carries out the same audit with a second,
independent input $V_3$ defined over $\Q$ alone, for which the rational
$2$-function condition is proved at \emph{every} prime, with no ramified
exception and no field extension. Both witnesses are retained: the one here
is smaller, the one there has no hypothesis that has to be read off a
discriminant. The subscripts distinguish them throughout; note in particular
that their weighted harmonic sums, $361/144$ and $361/16$, are different
numbers computed for different inputs, and neither is a misprint for the
other.

A different period-four counterexample is already documented by the same
repository that supplied the prior odd-prime proof \cite{priorcounter}.
Accordingly, the discovery of a problem in the general framing statement is
not claimed as new here. The spacing-three examples are compact independent
boundary witnesses. The criticism is specifically of the printed uniform
bounds in version 1. A single exceptional prime does not refute a weaker
assertion that permits an unspecified finite set of additional exceptional
primes, or one with a compensating scalar, and it says nothing by itself
about unrelated framing results. None of the proofs in this article uses the
disputed statements.

\section{Asymptotic expansion of the basket sequence}
\label{sec:asymptotic}

The leading asymptotic is recorded in the OEIS entry, attributed to
Kotesovec \cite{oeis348410}. Here we derive it together with two explicit
correction terms and an all-order coefficient formula. The method is the
standard saddle-point approach to large powers of generating functions
\cite[Chapter VIII]{flajolet}.

\subsection{The saddle and exponential growth}
For this section only, put
\[
 \Phi(x)=\frac1{(1-x)(1-x^2)}.
\]
The positive saddle $\tau\in(0,1)$ satisfies
\[
 \tau\frac{\Phi'(\tau)}{\Phi(\tau)}=1.
\]
Since $x\Phi'/\Phi=x/(1-x)+2x^2/(1-x^2)$, this is equivalent to
$4\tau^2+\tau-1=0$. Hence
\begin{equation}\label{eq:saddle}
 \tau=\frac{\sqrt{17}-1}{8},\qquad
 \rho=\frac{\tau}{\Phi(\tau)}
 =\frac{51\sqrt{17}-107}{512},\qquad
 \lambda=\rho^{-1}=\frac{107+51\sqrt{17}}{64}.
\end{equation}
Numerically, $\lambda=4.9574747954140732506\ldots$.
Define the logarithmic cumulants
\begin{equation}\label{eq:kappa}
 \kappa_j=\left.
 \left(x\frac{\mathrm d}{\mathrm dx}\right)^j\log\Phi(x)
 \right|_{x=\tau}.
\end{equation}
Then $\kappa_1=1$ and
\begin{align}\label{eq:kappas}
 \kappa_2&=\frac{51-5\sqrt{17}}{16},&
 \kappa_3&=\frac{435-69\sqrt{17}}{32},\notag\\
 \kappa_4&=\frac{10731-2125\sqrt{17}}{128},&
 \kappa_5&=\frac{91545-20175\sqrt{17}}{128},\notag\\
 \kappa_6&=\frac{4008567-929185\sqrt{17}}{512}.&&
\end{align}

\begin{theorem}[Two correction terms]\label{thm:asymptotic}
As $n\to\infty$ through positive integers,
\begin{equation}\label{eq:asymptotic}
 a(n)=\frac{\lambda^n}{\sqrt{2\pi\kappa_2n}}
 \left(1+\frac{c_1}{n}+\frac{c_2}{n^2}+O(n^{-3})\right),
\end{equation}
where
\begin{equation}\label{eq:c1c2}
 c_1=\frac{-1683+95\sqrt{17}}{9248},\qquad
 c_2=\frac{119821-28605\sqrt{17}}{5030912}.
\end{equation}
The leading amplitude is equivalently
\begin{equation}\label{eq:amplitude}
 \frac1{\sqrt{2\pi\kappa_2}}=
 \frac{\sqrt{3+5/\sqrt{17}}}{4\sqrt\pi}.
\end{equation}
\end{theorem}

\subsection{Derivation, including control of the remainder}
Cauchy's coefficient formula on the circle $|x|=\tau$ gives
\begin{equation}\label{eq:cauchy}
 a(n)=\frac{\lambda^n}{2\pi}
 \int_{-\pi}^{\pi}
 \exp\left(n\left\{
 \log\frac{\Phi(\tau e^{i\theta})}{\Phi(\tau)}-i\theta
 \right\}\right)\dd\theta.
\end{equation}
The coefficients of $\Phi$ are positive and its support contains $0$ and
$1$. Equality in the triangle inequality
$|\Phi(\tau e^{i\theta})|\le\Phi(\tau)$ therefore requires
$e^{i\theta}=1$. On every closed subinterval bounded away from zero, the
integrand is exponentially smaller than the central contribution.

Near zero, choose the analytic logarithm agreeing with the real logarithm
at $\tau$. The saddle condition cancels the linear term, and the exponent
inside braces is
\begin{equation}\label{eq:local-exponent}
 -\frac{\kappa_2\theta^2}{2}
 +\sum_{k\ge3}\frac{\kappa_k(i\theta)^k}{k!}.
\end{equation}
In particular, its real part is at most $-c\theta^2$ for sufficiently small
$|\theta|$ and some $c>0$.

Set $u=\sqrt{n\kappa_2}\,\theta$. The central integral becomes a Gaussian
integral multiplied by
\[
 \exp\left(\sum_{k\ge3}
 \frac{\kappa_k i^k u^k}
 {k!\kappa_2^{k/2}}n^{1-k/2}\right).
\]
Expanding this expression and integrating term by term is justified at any
fixed order by analytic Taylor remainder estimates and Gaussian domination.
For specificity, split a small fixed central interval at
$|\theta|=n^{-2/5}$. Outside that shrinking interval the bound
$e^{-cn\theta^2}$ is exponentially small in $n^{1/5}$. Inside it, Taylor
expansion to sufficiently high finite degree has a uniform remainder;
rescaling and integrating its polynomial bounds against a slightly wider
Gaussian gives the stated power remainder. Odd powers have zero symmetric
Gaussian moment. Carrying the expansion two orders further than the last
retained even power bounds the relative remainder by $O(n^{-3})$ in
\eqref{eq:asymptotic}. The same argument works to every fixed order.

Using
\[
 \frac1{\sqrt{2\pi}}\int_{\R}u^{2\ell}e^{-u^2/2}\dd u=(2\ell-1)!!,
\]
one obtains
\begin{align}\label{eq:cumulant-corrections}
 c_1={}&\frac{\kappa_4}{8\kappa_2^2}
       -\frac{5\kappa_3^2}{24\kappa_2^3},\notag\\
 c_2={}&-\frac{\kappa_6}{48\kappa_2^3}
 +\frac{7\kappa_3\kappa_5}{48\kappa_2^4}
 +\frac{35\kappa_4^2}{384\kappa_2^4}\notag\\
 &-\frac{35\kappa_3^2\kappa_4}{64\kappa_2^5}
 +\frac{385\kappa_3^4}{1152\kappa_2^6}.
\end{align}
Substitution of \eqref{eq:kappas} simplifies to \eqref{eq:c1c2}; elementary
rationalization gives \eqref{eq:amplitude}. This proves the theorem.

\subsection{An explicit all-order coefficient formula}
The same calculation gives a finite expression for every correction
coefficient. For $j\ge0$, sum over nonnegative integers
$m_3,\ldots,m_{2j+2}$ satisfying
\[
 \sum_{k=3}^{2j+2}(k-2)m_k=2j,
 \qquad K=\sum_{k=3}^{2j+2}km_k.
\]
The integer $K$ is even. Then
\begin{equation}\label{eq:allorder}
 c_j=\sum
 (-1)^{K/2}(K-1)!!\,\kappa_2^{-K/2}
 \prod_{k=3}^{2j+2}\frac{\kappa_k^{m_k}}{m_k!(k!)^{m_k}}.
\end{equation}
For $j=0$, the empty product and $(-1)!!$ are both $1$, giving $c_0=1$.
The constraint comes from collecting the exponent of $n^{-1/2}$ in the
expanded exponential; the double factorial comes from its Gaussian moment.
This is a direct finite enumeration formula, not a recurrence with unknown
initial correction terms.

Numerically, $c_1\approx-0.1396307272454933$ and
$c_2\approx0.0003736029529450764$. The cancellation in $c_2$ is substantial,
so exact symbolic simplification is preferable to low-precision arithmetic.
The supplied numerical checks use 90 decimal digits. For illustration,
the signed relative error of the approximation through $c_2$ is about
$-9.6264\cdot10^{-9}$ at $n=100$ and $-6.1254\cdot10^{-10}$ at $n=250$.
Those numerical comparisons are checks, not the proof of the remainder.

\section{Reproducible computation and proof audit}
\label{sec:verification}

\subsection{What the programs compute}
The archive separates exact integer checks from optional symbolic and
high-precision checks. The program \texttt{verify.py} uses only the Python
standard library and implements \eqref{eq:coefficient-recurrence}; it
cross-checks the target against \eqref{eq:positive-sum}, against
\eqref{eq:signed-sum}, and against a fourth, wholly independent mechanism,
namely repeated convolution of the single-pair basket series
$\lfloor j/2\rfloor+1$. It also tests
Lemmas~\ref{lem:gcoef}, \ref{lem:tail} and~\ref{lem:quadratic} directly, the
boundary family \eqref{eq:bn-boundary}, the rational $2$-function hypothesis
used in Appendix~\ref{app:audit}, the M\"obius denominators, and the
generating-function identities of Section~\ref{sec:gf} in integer arithmetic
to degree $60$. The program \texttt{binary\_checks.py} tests the complete
leading-defect formula, the harmonic-unit lemma, and both common-denominator
assertions. The program \texttt{symbolic\_checks.py} uses SymPy and mpmath
for elimination, Gaussian coefficient enumeration, and asymptotic
comparisons.

The recorded run produced the following exact checks. Each count is a finite
coverage statement, not an extrapolation to all indices.
\begin{center}
\small
\begin{tabular}{p{0.70\textwidth}r}
\toprule
Check & Cases\\
\midrule
Initial target terms, against the inspected OEIS listing &26\\
Target recurrence versus positive binomial sum &208\\
Target recurrence versus independent basket convolution &13\\
Signed parameter recurrence versus binomial sum &1,519\\
Target odd-prime congruences, strong form &489\\
Two-parameter odd-prime congruences, strong form &3,675\\
Odd-prime quadratic-tail tests &918\\
Logarithmic-derivative coefficient bounds &3,415\\
Factorial inequalities for the exponential tail &119,760\\
Generating-function coefficients, integer arithmetic &61\\
Original M\"obius/orbit checks &300\\
Cubic M\"obius denominator indices &100\\
Exact boundary witnesses &7\\
Period-three boundary family members &3\\
Rational $2$-function input checks for the audit &10,000\\
\midrule
Complete binary leading-defect formula &5,808\\
Binary first-lift congruence &1,936\\
Exact binary harmonic-unit normalization &400\\
Target denominator-12 and orbit divisibility &1,000\\
General denominator-24 / conditional denominator-12 &4,900\\
Binary sharpness and minimal-denominator witnesses &290\\
\bottomrule
\end{tabular}
\end{center}
All these checks passed, in a single run of the merged programs; no count
above is spliced from a different sweep. The binary leading-defect grid
covers $-5\le\alpha,\beta\le5$, odd $1\le m\le15$, and $2\le r\le7$, with
$m2^r\le2000$. The target odd-prime run uses odd primes at most $97$,
multipliers at most $12$ \emph{including} those divisible by the prime, and
indices at most $5000$; the largest index actually reached is $5000$ itself,
at $p=5$, $m=8$, $r=4$. Because
multipliers divisible by $p$ are no longer skipped, the run exercises the
strong exponent $3(r+v_p(m))-\eps_p$ of \eqref{eq:strong-m}, not merely the
coprime case. The general odd-prime grid uses parameters from $-3$ to $3$,
primes $3,5,7,11,13$, and indices at most $350$. Zero differences are
recorded as having infinite valuation. The derivative and quadratic checks
deliberately include signed and degenerate parameter pairs, for which some
prescribed bounds are zero and some coefficients vanish; those cases are
counted rather than silently discarded.

The symbolic run additionally verifies the logarithmic derivative, the
quartic resultant together with the constant quotient of that resultant by
the displayed quartic, its parametrization, its power series through degree
$20$, the three parity-weight polynomial identities used in
Lemma~\ref{lem:quadratic}, the saddle point, the growth constant, the leading
amplitude, and the Gaussian formulas for $c_1$ and $c_2$. These tests used
Python 3.13.5, SymPy 1.14.0, and mpmath 1.3.0 in the recorded environment;
the exact scripts are not tied to those library versions.

\subsection{Reproduction commands and output data}
From the archive directory, run
\begin{lstlisting}
python verify.py
python binary_checks.py
python -m pip install -r requirements-optional.txt
python symbolic_checks.py
latexmk -pdf -interaction=nonstopmode -halt-on-error article.tex
\end{lstlisting}
Only the third and fourth commands need optional Python packages. No network
access is used by any of the checking programs themselves, and no
floating-point arithmetic occurs in any congruence check: every such test is
performed in exact integer, rational, or modular arithmetic. The integer
checker accepts command-line coverage parameters; \texttt{--max-n 20000}
extends the original-sequence sweep, and \texttt{python verify.py --help}
lists the rest. In PowerShell the script path may also be written
\texttt{python .\textbackslash verify.py}.

The CSV files contain target terms through $1000$, the exact prime-power
valuation checks, the full binary parameter grid, primitive-orbit and
M\"obius values, the period-three boundary family, and asymptotic
comparisons. The reports preserve the actual check counts. There are no
randomized tests and no unrecorded dependence on an external computer-algebra
session.

\subsection{Source and artifact notes}
\label{sec:artifacts}
Sources were inspected on 19 September 2026. The framing theorem statements
of \cite{muller} were checked against the PDF of arXiv:2104.10754v1 --- page
3 for Theorem 1.1 and page 4 for Theorem 1.2, with Proposition 5.2(3) for
the framing coefficient formula --- and not only against the HTML rendering,
because the two can differ in the precise quantifiers. The prior proof
candidate and the prior period-four counterexample \cite{prior,priorcounter}
were read at a pinned immutable commit, whose full hash appears in the
bibliography. No database edit, no OEIS submission, and no communication
with any author was performed. No source PDFs, third-party repository trees,
or font files are bundled with this archive.

\subsection{Independent points of possible failure}
Several interfaces deserve explicit scrutiny in any independent review.
First, a recurrence in the coefficient index must not be mistaken for a
recurrence in the diagonal index. Second, the odd-prime exponential has a
vanishing linear term, but the binary one does not; the latter is handled
by reflection, not by reusing the former proof. Third, the $1/2$ in formal
integration by parts at the binary prime is retained throughout and is
responsible for one of the losses. Fourth, the Möbius grouping is performed
at the exact prime-power exponent of $n$, not merely at a prime dividing
$n$. Finally, the finer binary valuation table is not used to justify any
of the uniform bounds.

The article uses only elementary formal series, local rational arithmetic,
coefficient extraction, and the stated saddle-point analysis. It does not
invoke the disputed periodic framing bound or any unverified asymptotic
recurrence from an OEIS comment.

\section{Further questions after the binary completion}
\label{sec:frontiers}

The uniform small-prime completion is proved, but it does not determine
all exceptional extra divisibility. For the basket sequence, the observed
valuations at $n=2^r$ are
\begin{center}
\begin{tabular}{r r r r r r r r r}
\toprule
$r$&1&2&3&4&5&6&7&8\\
\midrule
$v_2(a(2^r)-a(2^{r-1}))$&2&4&8&13&16&19&22&25\\
\bottomrule
\end{tabular}
\end{center}
The recorded exact computation continues this pattern through $r=15$.

\begin{conjecture}[A finer target valuation, not used above]\label{conj:finer}
For every $r\ge4$,
\[
 v_2\bigl(a(2^r)-a(2^{r-1})\bigr)=3r+1.
\]
\end{conjecture}

This is not a consequence of Theorem~\ref{thm:binary}, whose leading residue
vanishes for $(\alpha,\beta)=(2,1)$. Proving it would require following the
reflection calculation several binary digits further, including the next
residues of $K(t)$ and additional exponential terms. The statement for a
fixed odd multiplier is more varied: at $m=5$ the values from $r=4$ onward
in the checked range are $3r+3$, while at $m=7$ they are $3r+2$.

A second question is to characterize the periodic logarithmic coefficient
sequences for which the cubic tail estimate holds. Formula
\eqref{eq:obstruction} gives an immediate first-level obstruction, but
vanishing at the first level alone does not establish all higher levels.
The spacing-three example shows why a condition stronger than periodicity
and Frobenius invariance is needed: any classification of the coefficient
patterns satisfying the criterion \eqref{eq:criterion} must accommodate the
period-three obstruction rather than assume it away. A useful classification
would separate identities forced by symmetry from congruences that depend on
the prime.

Closely related, and equally open, is the question of strictness. Neither
Theorem~\ref{thm:odd} nor Theorem~\ref{thm:binary} classifies the parameters,
primes and levels at which its lower bound is attained rather than exceeded.
Table~\ref{tab:valuations} shows equality throughout its odd range, and
Theorem~\ref{thm:binary} settles the binary case for odd $\alpha\beta$ and
$m=1$; between these the picture is empirical.

A third question is to refine the common denominator for a fixed parameter
pair. The bounds $24$ and $12$ are optimal uniformly over the stated classes,
but special pairs can have smaller denominators. The exact first binary
residue already detects one such improvement. An analogous explicit
ternary residue could identify further parameter-dependent reductions.

These questions are distinct from the completed claims. In particular,
the remaining uncertainty about the finer valuation pattern does not limit
the proof of the universal denominator-12 theorem.

\section{Conclusion}

Two theorems have been proved, by two different mechanisms. At an odd prime,
the parity of the logarithmic coefficients survives multiplication by $p$,
which produces an exact factorization; the whole difficulty is then
concentrated in one weighted reciprocal-square sum, and
Lemmas~\ref{lem:squares} and~\ref{lem:quadratic} show that it supplies the
third power of $p$, with a single unavoidable loss at $p=3$ traceable to the
valuation of the constant $3$. At $p=2$ that parity argument is unavailable,
and the binary completion instead converts a linear exponential contribution
into a quadratic one by reflection and then extracts an extra index factor by
formal integration by parts. An exact odd-harmonic normalization identifies
the first possible binary bit, not just a lower bound. Priority for the
odd-prime statement belongs to the prior proof candidate located in the
source audit; the binary results were not found in the checked sources.

For A348410 the completed all-prime estimate is naturally encoded by the
single integer $12$. It clears every denominator of the cubic M\"obius
transform, is forced by the values at $n=3$ and $n=4$, and is already
attained in full at the single index $n=12$. The resulting primitive
cyclic-word divisibility gives a combinatorial expression of the same
arithmetic fact.

The negative results are worth keeping in view beside the positive ones.
Exact invariance of the logarithmic coefficients under index multiplication
is not enough; the quadratic estimate \eqref{eq:criterion} is a separate
hypothesis, and an infinite period-three family violates it at $p=5$. The
same family refutes the explicit conclusion of a printed general framing
statement, which is why no proof here cites that statement.

The archive contains the full proofs and reproducible checks; priority beyond
the audited sources and external validation of the new binary argument remain
matters for further review.

\appendix
\section{A minimal exact evaluator}
The following code is the computational core; the full scripts add input
validation, independent formula checks, report generation, and CSV output.
\begin{lstlisting}[language=Python]
def coefficient_diagonal(n: int, alpha: int = 2,
                         beta: int = 1) -> int:
    if n < 0:
        raise ValueError("n must be nonnegative")
    f_minus_two, f_minus_one = 0, 1
    for k in range(1, n + 1):
        numerator = n * (alpha - beta) * f_minus_one
        numerator += (n * (alpha + beta) + k - 2) * f_minus_two
        current, remainder = divmod(numerator, k)
        if remainder:
            raise ArithmeticError("coefficient division was not exact")
        f_minus_two, f_minus_one = f_minus_one, current
    return f_minus_one
\end{lstlisting}
The running names refer to the two previously available coefficients at the
start of each iteration. At $n=0$ the function correctly returns $1$.

\section{Notation and dependency map}
\begin{center}
\small
\begin{tabular}{p{0.20\textwidth}p{0.73\textwidth}}
\toprule
Symbol & Meaning\\
\midrule
$\A(n)$ & Integer-parameter coefficient diagonal in \eqref{eq:family}.\\
$a(n)$ & The specialization $A_{2,1}(n)$, OEIS A348410.\\
$A^{(d)}(n)$ & The change-of-step family of Corollary~\ref{cor:step}.\\
$\eps_p$ & The ternary loss \eqref{eq:epsilon}: $1$ at $p=3$, else $0$.\\
$R_p$ & Rationals with denominator prime to $p$.\\
$c_j$ & Logarithmic-derivative coefficients, \eqref{eq:phi-log}.\\
$g_j$ & Coefficients of a power $G^M$, Lemma~\ref{lem:gcoef}.\\
$D_p$ & Reduced odd-prime logarithm in \eqref{eq:Dp}.\\
$H_p(s)$ & Reduced reciprocal-square sum; not the generating function.\\
$T_p(s)$ & Quadratic tail coefficient $[x^s]D_p^2$, \eqref{eq:Tvaluation}.\\
$h(T)$ & Sum of reciprocals of the positive odd integers less than $T$.\\
$U,K,V$ & Binary series in \eqref{eq:U} and \eqref{eq:K}; $V$ is local to
Section~\ref{sec:binary}.\\
$b_t(n)$ & Period-three boundary family \eqref{eq:bn-boundary};
$b=b_1$.\\
$V_1,V_3$ & The two framing audit inputs, \eqref{eq:V1-input} and
\eqref{eq:V3-input}.\\
$B(n)$ & Cubic M\"obius transform $n^{-3}\sum_{d\mid n}\mu(d)a(n/d)$.\\
$P_n$ & Primitive cyclic-word count, equal to $n^2B(n)$.\\
$H(z)$ & Ordinary generating function of $a(n)$.\\
$\kappa_j$ & Saddle cumulants in \eqref{eq:kappa}.\\
\bottomrule
\end{tabular}
\end{center}
The proof dependencies are short. Lemmas~\ref{lem:gcoef}--\ref{lem:quadratic}
imply Theorem~\ref{thm:odd}, and hence Corollary~\ref{cor:strong-m}.
Lemma~\ref{lem:binary-H}, reflection, and
\eqref{eq:ibp} imply Theorem~\ref{thm:binary}. These two theorems plus the
first binary lift imply Theorem~\ref{thm:allprime}. Grouping M\"obius divisors
then proves Theorem~\ref{thm:denom}; periodic word decomposition gives
Corollary~\ref{cor:necklace}. The generating-function and asymptotic sections
are independent checks and complementary structure, not premises of the
arithmetic proofs. Appendix~\ref{app:audit} is a negative result about an
external statement and is used by nothing above it.

\section{Audit of a broader framing claim}
\label{app:audit}

An apparently relevant external source is M\"uller's preprint
\cite{muller}, arXiv:2104.10754v1. Its general statements, if correct as
printed, would imply a cubic prime-power congruence for a much wider class
than Theorem~\ref{thm:odd} covers, and it would be tempting to cite them as
a black box in place of Section~\ref{sec:odd}. This appendix records why we
do not. Nothing in the article depends on it.

In the coefficient convention of that preprint, a \emph{rational
$2$-function} over $\Q$ is a rational series $V(x)=\sum_{j\ge1}c_jx^j$ with
integral coefficients satisfying
\begin{equation}\label{eq:two-function}
 c_{mp^r}\equiv c_{mp^{r-1}}\pmod{p^{2r}}
 \quad\text{for every prime }p\text{ and all }m,r\ge1.
\end{equation}
For framing parameter $1$, Proposition 5.2(3) there gives the transformed
coefficients
\begin{equation}\label{eq:framing}
 \mathrm{fr}(n)=[x^n]\exp\Bigl(n\sum_{j\ge1}\frac{c_j}{j}x^j\Bigr),
\end{equation}
and the explicit $p\ge5$ conclusion of its Theorem 1.1 asserts cubic
prime-power congruences for $\mathrm{fr}$ whenever $p$ does not divide the
coefficient period.

Section~\ref{sec:litscope} already exhibited one input satisfying the printed
hypotheses and violating that conclusion, namely $V_1$ of
\eqref{eq:V1-input}, read over $K=\Q(\sqrt{-3})$. The witness below is
independent of it and is kept because it buys something the first does not:
its rational $2$-function property is proved at every prime by a finite case
split, with no ramified-prime exception and no field extension anywhere in
the argument.

\subsection{An input that satisfies the condition at every prime}
Take
\begin{equation}\label{eq:V3-input}
 V_3(x)=\frac{3x}{1-x}+\frac{9x^3}{1-x^3},
 \qquad c_j=3+9\cdot\mathbf1_{3\mid j}.
\end{equation}
The coefficient sequence is integral and has period $3$. We verify
\eqref{eq:two-function} for all indices, not merely on a finite test range.
Let $p$ be a prime and $m,r\ge1$.
\begin{itemize}[leftmargin=1.4em,itemsep=2pt,topsep=2pt]
\item If $p\ne3$, then multiplication by $p$ preserves divisibility by $3$,
so $3\mid mp^r$ if and only if $3\mid mp^{r-1}$, and the two coefficients are
\emph{equal}.
\item If $p=3$ and $r\ge2$, both $m3^r$ and $m3^{r-1}$ are divisible by $3$,
so again the two coefficients are equal.
\item If $p=3$ and $r=1$, the difference $c_{3m}-c_m$ is $0$ when $3\mid m$
and $9$ otherwise. In both cases it is divisible by $3^2=p^{2r}$.
\end{itemize}
Hence $V_3$ is a rational $2$-function over $\Q$ in the stated convention,
with no ramified-prime exception and no passage to a number field.

\subsection{The refutation arithmetic}
Formal integration of \eqref{eq:V3-input} gives
\[
 \sum_{j\ge1}\frac{c_j}{j}x^j=-3\log(1-x)-3\log(1-x^3),
\]
so that \eqref{eq:framing} becomes, in the notation of
\eqref{eq:bn-boundary},
\[
 \mathrm{fr}(n)=[x^n](1-x)^{-3n}(1-x^3)^{-3n}=b_3(n).
\]
At $n=1$ and $n=5$,
\begin{align*}
 \mathrm{fr}(1)&=3,\\
 \mathrm{fr}(5)&=\binom{19}{5}+15\binom{16}{2}
       =11\,628+1\,800=13\,428,\\
 \mathrm{fr}(5)-\mathrm{fr}(1)&=13\,425=25\cdot537\equiv50\pmod{125}.
\end{align*}
The prime $5$ is unramified over $\Q$ and does not divide the period $3$, so
Theorem 1.1 as printed would require divisibility by $5^3=125$. It fails.
This is the member $t=3$ of the family \eqref{eq:period3-family}, whose exact
valuation there is $2$.

\subsection{The weighted harmonic assertion fails on the same input}
Theorem 1.2 of the same preprint, repeated there as Theorem 6.2, would
require the period-three weighted sum below to vanish modulo $5$ at
$p=n=5$. For the very same coefficients \eqref{eq:V3-input}, exact rational
arithmetic gives
\begin{equation}\label{eq:audit-harmonic}
 \sum_{k=1}^{4}\frac{c_{5-k}c_k}{k^2}
 =9+\frac{36}{4}+\frac{36}{9}+\frac9{16}
 =\frac{361}{16}\equiv1\pmod5,
\end{equation}
and the corresponding quadratic coefficient is $33/2$, which likewise does
not vanish modulo $5$.

The value $361/16$ here and the value $361/144$ of
\eqref{eq:weighted-counterexample} are computed for \emph{different} inputs
--- $V_3$ with $c_j=3+9\cdot\mathbf1_{3\mid j}$, and $V_1$ with
$c_j=1+3\cdot\mathbf1_{3\mid j}$ respectively --- and their residues modulo
$5$ differ accordingly, $1$ against $4$. Neither is a misprint for the other.
Both are checked exactly by the supplied programs.

\subsection{Scope of this appendix}
The scope of the observation matters, and it is narrow. A failure at one
prime refutes the cited explicit all-eligible-primes conclusion, as printed
in the inspected version. It does not by itself refute a weaker statement
that allows an unspecified finite set of exceptional primes, or a
compensating scalar, and it says nothing about the author's intentions, about
an uninspected later revision, or about unrelated framing results. A
period-four counterexample is in any case already documented elsewhere
\cite{priorcounter}, so no novelty is claimed for the observation itself.

This audit is included because, without it, it would be tempting to cite the
framing statement as a black box. The main proof does not do that. Its
logarithmic expansion and coefficient estimate are in the same mathematical
tradition, but the reciprocal-square statement it needs is proved directly
for parity weights in Lemmas~\ref{lem:squares} and~\ref{lem:quadratic}. That
replacement is sufficient for A348410 and its integer-parameter family, and
it does not inherit the false arbitrary-period assertion.

\begin{thebibliography}{99}
\bibitem{oeis348410}
OEIS Foundation Inc., \emph{A348410: Number of nonnegative integer solutions
to $n=\sum_{i=1}^n(a_i+b_i)$, with $b_i$ even}. Includes P.~Bala's
21 February 2022 coefficient formulas and supercongruence conjecture,
and V.~Kotesovec's leading asymptotic.
\url{https://oeis.org/A348410}. Inspected 19 September 2026.

\bibitem{oeis352373}
OEIS Foundation Inc., \emph{A352373}. Related parameterized
coefficient-power congruence.
\url{https://oeis.org/A352373}. Inspected 19 September 2026.

\bibitem{prior}
GitHub repository \texttt{rbajaj5/a183068-supercongruence},
\emph{A cubic tower for a two-parameter coefficient family},
\texttt{related-results/CoefficientFramingCubicTower.md}.
Public proof candidate; pinned repository snapshot
\texttt{3085b46fdbce} (full hash in the linked URL).
\href{https://github.com/rbajaj5/a183068-supercongruence/blob/3085b46fdbce945d156d2b5b8b9d1a66627b4375/related-results/CoefficientFramingCubicTower.md}{Permanent link to the inspected note}.
Inspected 19 September 2026. The source states no binary assertion.

\bibitem{priorcounter}
GitHub repository \texttt{rbajaj5/a183068-supercongruence},
\emph{Counterexample to the rational-framing theorem},
\texttt{related-results/RationalFramingCounterexample.md}, same pinned
snapshot as \cite{prior}.
\href{https://github.com/rbajaj5/a183068-supercongruence/blob/3085b46fdbce945d156d2b5b8b9d1a66627b4375/related-results/RationalFramingCounterexample.md}{Permanent link to the period-four counterexample}.
Inspected 19 September 2026; also cross-checked against the source's
result index and coefficient-framing discussion.

\bibitem{niu}
T.~Niu, \emph{An explicit algebraic generating function for OEIS A348410},
arXiv:2605.16553, 2026.
\url{https://arxiv.org/abs/2605.16553}.

\bibitem{prodinger}
H.~Prodinger, \emph{The generating function of A348410 in OEIS using the
diagonal method and another sequence (A001008) from OEIS},
arXiv:2605.21255, version 2, 23 May 2026.
\url{https://arxiv.org/abs/2605.21255}.

\bibitem{muller}
L.~Felipe M\"uller, \emph{Wolstenholme Type Congruences and Framing of Rational
2-Functions}, arXiv:2104.10754, version 1, 21 April 2021.
Cited loci: Theorem 1.1, Theorem 1.2 (repeated as Theorem 6.2), and
Proposition 5.2(3), the framing coefficient formula.
\url{https://arxiv.org/abs/2104.10754}.
Statements checked against the PDF, pages 3 and 4, and not only the HTML
rendering; inspected 19 September 2026.
The criticism in Section~\ref{sec:boundary} and Appendix~\ref{app:audit}
concerns the precise printed bounds of Theorems 1.1, 1.2, and 6.2 in this
version, not all framing theory.

\bibitem{flajolet}
P.~Flajolet and R.~Sedgewick, \emph{Analytic Combinatorics}, Cambridge
University Press, 2009, especially Chapter VIII.
Authors' book site: \url{https://ac.cs.princeton.edu/home/}.
\end{thebibliography}
\end{document}
